\clearpage
\thispagestyle{empty}\SourceLeaf{201}
\pdfbookmark[0]{Allgemeine Lehrsätze über Anziehungs- und Abstoßungskräfte, die umgekehrt proportional zum Quadrat der Entfernung wirken}{allgemeine-lehrsaetze-title}
\begin{center}
\vspace*{0.6in}
{\Large ALLGEMEINE LEHRSÄTZE\par}
\bigskip
ÜBER ANZIEHUNGS- UND ABSTOSSUNGSKRÄFTE,\par
\bigskip
DIE UMGEKEHRT PROPORTIONAL ZUM\par
\bigskip
{\Large QUADRAT DER ENTFERNUNG WIRKEN\par}
\vspace{0.5in}
VON\par
\bigskip
{\large CARL FRIEDRICH GAUSS.\par}
\vfill
\rule{0.85\textwidth}{0.3pt}\par
Resultate aus den Beobachtungen des magnetischen Vereins im Jahre 1839.\par
Herausgegeben von \textsc{Gauss} und \textsc{Weber}. Leipzig 1840.\par
\rule{0.85\textwidth}{0.3pt}
\end{center}
\clearpage
\SourceLeaf{203}
\fancyhead[L]{\small Allgemeine Lehrsätze über Anziehungs- und Abstoßungskräfte}
\renewcommand{\Article}[1]{\par\medskip\Needspace{4\baselineskip}\pdfbookmark[1]{Art. #1}{lehrsaetze-art-#1}\centerline{\textbf{#1.}}\nobreak\smallskip}
\pdfbookmark[0]{Allgemeine Lehrsätze über Anziehungs- und Abstoßungskräfte}{lehrsaetze-body}
\begin{center}
ALLGEMEINE LEHRSÄTZE\\
ÜBER ANZIEHUNGS- UND ABSTOSSUNGSKRÄFTE,\\
DIE UMGEKEHRT PROPORTIONAL ZUM\\
QUADRAT DER ENTFERNUNG WIRKEN.
\end{center}
\Article{1}
In der Natur begegnen uns mancherlei Erscheinungen, die wir durch die Annahme erklären, dass die kleinsten Teilchen der Stoffe Kräfte aufeinander ausüben, die umgekehrt proportional zu den Quadraten ihrer gegenseitigen Entfernungen sind.

Hierher gehört vor allem die allgemeine Gravitation Durch sie übt jedes wägbare Molekül $\mu$ auf ein anderes $\mu'$ eine bewegende Kraft aus. Setzt man die Entfernung $=r$, so wird diese Kraft durch $\frac{\mu\mu'}{rr}$ ausgedrückt und wirkt darauf hin, dass sich die Moleküle entlang der geraden Verbindungslinie einander nähern.

Nimmt man zur Erklärung der magnetischen Erscheinungen zwei magnetische Flüssigkeiten an, von denen die eine als positive Größe und die andere als negative gilt, so üben zwei solche Elemente $\mu$, $\mu'$ ebenfalls eine bewegende Kraft aufeinander aus. Sie wird durch $\frac{\mu\mu'}{rr}$ gemessen und wirkt entlang der geraden Verbindungslinie als Abstoßung, wenn $\mu$, $\mu'$ gleichartig sind, und als Anziehung, wenn sie verschiedenartig sind.

Ganz Ähnliches gilt für die gegenseitige Wirkung der Teilchen der elektrischen Flüssigkeiten.

\SourceLeaf{204}
Das lineare Element $\mathrm{d}s$ eines galvanischen Stroms übt auf ein Element $\mu$ der magnetischen Flüssigkeit (wenn man diese gelten lässt) ebenfalls eine bewegende Kraft aus, die umgekehrt proportional zum Quadrat der Entfernung $r$ ist. aber hier kommt ein ganz anderer Umstand hinzu. Die Kraft wirkt senkrecht zu der Ebene, die durch $\mu$ und die Richtung von $\mathrm{d}s$ gelegt wird, statt entlang der geraden Verbindungslinie, und ihre Stärke hängt außer von der Entfernung auch von dem Winkel ab, den $r$ mit der Richtung von $\mathrm{d}s$ bildet. Nennt man diesen Winkel $\theta$, so misst $\frac{\sin\theta.\mu\mathrm{d}s}{rr}$ die bewegende Kraft, die $\mathrm{d}s$ auf $\mu$ ausübt, und $\mu$ übt auf das Stromelement $\mathrm{d}s$ oder dessen wägbaren Träger eine ebenso große Kraft aus, deren Richtung parallel zur Richtung der ersten Kraft und ihr entgegengesetzt ist.

Nimmt man mit \textsc{Ampère} an, dass zwei Elemente $\mathrm{d}s$, $\mathrm{d}s'$ galvanischer Ströme entlang ihrer geraden Verbindungslinie anziehend oder abstoßend aufeinander wirken, so zwingen uns die Erscheinungen, auch diese Kraft als umgekehrt proportional zum Quadrat der Entfernung anzusetzen, und verlangen zugleich eine etwas kompliziertere Abhängigkeit von der Richtung der Stromelemente.

In dieser Abhandlung werden wir uns auf die ersten drei Fälle beschränken, also auf Kräfte, die entlang der Geraden zwischen dem wirkenden Element und dem Element, auf das es wirkt, auftreten und schlechthin umgekehrt proportional zum Quadrat der Entfernung sind, obwohl sich mehrere Lehrsätze mit kleinen Änderungen auch auf die anderen Fälle anwenden lassen. Deren ausführliche Behandlung muss einer anderen Abhandlung vorbehalten bleiben.

\Article{2}
Wir bezeichnen mit $a$, $b$, $c$ die rechtwinkligen Koordinaten eines materiellen Punktes, von dem eine abstoßende oder anziehende Kraft ausgeht. Die beschleunigende Kraft selbst in einem beliebigen Punkt $O$ mit den Koordinaten $x$, $y$, $z$ bezeichnen wir mit
\[
\frac{\mu}{(a-x)^2+(b-y)^2+(c-z)^2}=\frac{\mu}{rr}
\]
Dabei steht $\mu$ im ersten Fall des vorhergehenden Artikels für die im ersten Punkt
\SourceLeaf{205}
vorhandene wägbare Materie, im zweiten und dritten für die Menge der magnetischen oder elektrischen Flüssigkeit. Zerlegt man diese Kraft parallel zu den drei Koordinatenachsen, so ergeben sich die Komponenten
\[
\frac{\varepsilon\mu(a-x)}{r^3},\qquad \frac{\varepsilon\mu(b-y)}{r^3},\qquad \frac{\varepsilon\mu(c-z)}{r^3}
\]
wobei $\varepsilon=+1$ oder $=-1$ sein soll, je nachdem, ob die Kraft anziehend oder abstoßend wirkt. Das ergibt sich aus der Beschaffenheit dessen, was die Kraft ausübt, und dessen, worauf sie wirkt. Diese Komponenten lassen sich als partielle Ableitungen darstellen
\[
\frac{\mathrm{d}\frac{\varepsilon\mu}{r}}{\mathrm{d}x},\qquad \frac{\mathrm{d}\frac{\varepsilon\mu}{r}}{\mathrm{d}y},\qquad \frac{\mathrm{d}\frac{\varepsilon\mu}{r}}{\mathrm{d}z}
\]
Wirken also auf denselben Punkt $O$ mehrere Kraftquellen $\mu^0$, $\mu'$, $\mu''$ usw. aus den Entfernungen $r^0$, $r'$, $r''$ usw., und setzt man
\[
\frac{\mu^0}{r^0}+\frac{\mu'}{r'}+\frac{\mu''}{r''}+\text{ u. s. f.}=\Sigma\frac{\mu}{r}=V
\]
so werden die Komponenten der gesamten in $O$ wirkenden Kraft durch
\[
\frac{\varepsilon\mathrm{d}V}{\mathrm{d}x},\qquad \frac{\varepsilon\mathrm{d}V}{\mathrm{d}y},\qquad \frac{\varepsilon\mathrm{d}V}{\mathrm{d}z}
\]
dargestellt.

Wenn die Kraftquellen eine Linie, eine Fläche oder ein räumliches Gebiet stetig ausfüllen, statt aus diskreten Punkten zu wirken, tritt an die Stelle der Summation $\Sigma$ eine einfache, doppelte oder dreifache Integration. In der Natur selbst kommt nur der letzte Fall vor. Da man ihn aber unter gewissen Einschränkungen oft durch gedachte Kraftquellen ersetzen kann, die in Punkten konzentriert oder auf Linien oder Flächen stetig verteilt sind, werden wir auch diese Fälle in unsere Untersuchung einbeziehen. Dabei können wir ohne Bedenken von \textit{Massen} sprechen, die auf einer Fläche oder einer Linie verteilt oder in einem Punkt konzentriert sind. Der Ausdruck Masse bezeichnet hier lediglich das, wovon die Anziehungs- oder Abstoßungskräfte ausgehen sollen.

\Article{3}
Wir bezeichnen also für jeden Punkt im Raum mit $x$, $y$, $z$ seine rechtwinkligen Koordinaten und mit $V$ die Summe aller wirkenden Massenteilchen, je\SourceLeaf{206}des geteilt durch seine Entfernung von diesem Punkt. Je nach den Bedingungen der jeweiligen Untersuchung können dabei negative Massenteilchen ausgeschlossen oder zugelassen werden. Damit wird $V$ eine Funktion von $x$, $y$, $z$, und die Erforschung der besonderen Eigenschaften dieser Funktion wird der Schlüssel zur Theorie der Anziehungs- oder Abstoßungskräfte selbst sein. Um die dazu nötigen Untersuchungen leichter handhaben zu können, erlauben wir uns, diesem $V$ eine eigene Bezeichnung zu geben und die Größe das \textit{Potential} der Massen zu nennen, auf die sie sich bezieht. Für unsere jetzige Untersuchung reicht diese engere Begriffsbestimmung aus. Im weiteren Sinn könnte man sowohl bei anderen Anziehungsgesetzen als dem Gesetz der umgekehrten Proportionalität zum Quadrat der Entfernung als auch im vierten in Artikel 1 erwähnten Fall unter Potential die Funktion von $x$, $y$, $z$ verstehen, deren partielle Ableitungen die Komponenten der erzeugten Kraft angeben.

Bezeichnen wir die gesamte im Punkt $x$, $y$, $z$ wirkende Kraft mit $p$ und die Winkel, die ihre Richtung mit den drei Koordinatenachsen bildet, mit $\alpha$, $\beta$, $\gamma$, so sind die drei Komponenten
\[
p\cos\alpha=\varepsilon\frac{\mathrm{d}V}{\mathrm{d}x},\qquad p\cos\beta=\varepsilon\frac{\mathrm{d}V}{\mathrm{d}y},\qquad p\cos\gamma=\varepsilon\frac{\mathrm{d}V}{\mathrm{d}z}
\]
und
\[
p=\BareRoot{((\frac{\mathrm{d}V}{\mathrm{d}x})^2+(\frac{\mathrm{d}V}{\mathrm{d}y})^2+(\frac{\mathrm{d}V}{\mathrm{d}z})^2)}
\]
\Article{4}
Ist $\mathrm{d}s$ das Element einer beliebigen geraden oder gekrümmten Linie, so sind $\frac{\mathrm{d}x}{\mathrm{d}s}$, $\frac{\mathrm{d}y}{\mathrm{d}s}$, $\frac{\mathrm{d}z}{\mathrm{d}s}$ die Kosinus der Winkel, die dieses Element mit den Koordinatenachsen bildet. Bezeichnet also $\theta$ den Winkel zwischen der Richtung des Elements und der Richtung der dort wirkenden resultierenden Kraft, so gilt
\[
\cos\theta=\frac{\mathrm{d}x}{\mathrm{d}s}.\cos\alpha+\frac{\mathrm{d}y}{\mathrm{d}s}.\cos\beta+\frac{\mathrm{d}z}{\mathrm{d}s}.\cos\gamma
\]
Die auf die Richtung von $\mathrm{d}s$ projizierte Kraft ist folglich
\[
p\cos\theta=\varepsilon(\frac{\mathrm{d}V}{\mathrm{d}x}.\frac{\mathrm{d}x}{\mathrm{d}s}+\frac{\mathrm{d}V}{\mathrm{d}y}.\frac{\mathrm{d}y}{\mathrm{d}s}+\frac{\mathrm{d}V}{\mathrm{d}z}.\frac{\mathrm{d}z}{\mathrm{d}s})=\frac{\varepsilon\mathrm{d}V}{\mathrm{d}s}
\]
Legen wir durch alle Punkte, an denen das Potential $V$ einen konstanten Wert hat, eine Fläche, so trennt sie im Allgemeinen die Teile des
\SourceLeaf{207}
Raumes, in denen $V$ kleiner ist als dieser Wert, von denen, in denen $V$ größer ist. Liegt die Linie $s$ in dieser Fläche oder verläuft wenigstens ihr Element $\mathrm{d}s$ tangential zu dieser Fläche, so gilt $\frac{\mathrm{d}V}{\mathrm{d}s}=0$. Wenn sich die Anteile der gesamten Kraft an dieser Stelle gegenseitig aufheben, wenn also $p=0$ wird, kann von einer Richtung der Kraft keine Rede mehr sein. Andernfalls muss $\cos\theta=0$ sein. Daraus schließen wir, dass die resultierende Kraft in jedem Punkt einer solchen Fläche senkrecht zu ihr steht. Sie weist auf die Seite des Raumes, auf der die größeren Werte von $V$ an die Fläche angrenzen, wenn $\varepsilon=+1$ ist, und auf die entgegengesetzte Seite, wenn $\varepsilon=-1$ ist. Wir nennen eine solche Fläche eine \textit{Gleichgewichtsfläche}. Da sich durch jeden Punkt eine solche Fläche legen lässt, trifft die Linie $s$ in jedem ihrer Punkte eine andere Gleichgewichtsfläche, sofern sie nicht ganz in einer einzigen solchen Fläche liegt. Schneidet $s$ alle Gleichgewichtsflächen im rechten Winkel, so gibt eine Tangente an diese Linie überall die Richtung der Kraft an und $\frac{\mathrm{d}V}{\mathrm{d}s}$ ihre Stärke.

Das Integral $\int p\cos\theta.\mathrm{d}s$ über ein beliebiges Stück der Linie $s$ ist offenbar $=\varepsilon(V'-V^0)$, wobei $V^0$, $V'$ die Werte des Potentials am Anfangs- und am Endpunkt bezeichnen. Ist $s$ also eine geschlossene Linie, so ist dieses Integral über die ganze Linie $=0$.

\Article{5}
Es ist klar, dass das Potential in jedem Punkt des Raumes, der \textit{außerhalb} aller anziehenden oder abstoßenden Teilchen liegt, einen angebbaren Wert haben muss. Dasselbe gilt für seine Ableitungen erster wie höherer Ordnung, da sie unter dieser Voraussetzung ebenfalls die Form von Summen angebbarer Teile oder von Integralen über Differentiale annehmen, deren Koeffizienten durchweg angebbare Werte haben. So erhalten wir
\[
\begin{array}{ll}
\displaystyle\frac{\mathrm{d}V}{\mathrm{d}x}=\Sigma\frac{(a-x)\mu}{r^3},&\displaystyle\frac{\mathrm{d}\mathrm{d}V}{\mathrm{d}x^2}=\Sigma(\frac{3(a-x)^2}{r^5}-\frac{1}{r^3})\mu\\[6pt]
\displaystyle\frac{\mathrm{d}V}{\mathrm{d}y}=\Sigma\frac{(b-y)\mu}{r^3},&\displaystyle\frac{\mathrm{d}\mathrm{d}V}{\mathrm{d}y^2}=\Sigma(\frac{3(b-y)^2}{r^5}-\frac{1}{r^3})\mu\\[6pt]
\displaystyle\frac{\mathrm{d}V}{\mathrm{d}z}=\Sigma\frac{(c-z)\mu}{r^3},&\displaystyle\frac{\mathrm{d}\mathrm{d}V}{\mathrm{d}z^2}=\Sigma(\frac{3(c-z)^2}{r^5}-\frac{1}{r^3})\mu
\end{array}
\]
\SourceLeaf{208}
Die bekannte Gleichung
\[
\frac{\mathrm{d}\mathrm{d}V}{\mathrm{d}x^2}+\frac{\mathrm{d}\mathrm{d}V}{\mathrm{d}y^2}+\frac{\mathrm{d}\mathrm{d}V}{\mathrm{d}z^2}=0
\]
gilt also für alle Punkte des Raumes, die außerhalb der wirkenden Massen liegen.

\Article{6}
Unter den verschiedenen Fällen, in denen es auf den Wert des Potentials $V$ oder seiner Ableitungen an einem Punkt ankommt, der nicht außerhalb der wirkenden Massen liegt, wollen wir zuerst den in der Natur vorkommenden Fall betrachten. Dabei füllen die Massen ein bestimmtes räumliches Gebiet mit gleichförmiger oder ungleichförmiger, aber überall endlicher Dichte aus.

Es sei $t$ das gesamte Gebiet, das Masse enthält, und $\mathrm{d}t$ ein unendlich kleines Element davon mit den Koordinaten $a$, $b$, $c$ und dem Massenelement $k\mathrm{d}t$. Ferner sei $V$ das Potential im Punkt $O$ mit den Koordinaten $x$, $y$, $z$. Die Entfernung von jenem Element ist also
\[
\BareRoot{((a-x)^2+(b-y)^2+(c-z)^2)}=r
\]
Folglich gilt
\[
V=\int\frac{k\mathrm{d}t}{r}
\]
wobei über das ganze Gebiet $t$ zu integrieren ist, was eine dreifache Integration bedeutet. Man sieht leicht, dass sich die Integration auch dann tatsächlich ausführen lässt, wenn $O$ innerhalb des Gebietes liegt, obwohl $\frac{1}{r}$ dann für die unendlich nahe bei $O$ liegenden Elemente unendlich groß wird. Führt man nämlich statt $a$, $b$, $c$ Polarkoordinaten ein, indem man
\[
a=x+r\cos u,\qquad b=y+r\sin u\cos\lambda,\qquad c=z+r\sin u\sin\lambda
\]
setzt, so ergibt sich $\mathrm{d}t=rr\sin u.\mathrm{d}u.\mathrm{d}\lambda.\mathrm{d}r$ und damit
\[
V=\int\!\int\!\int kr\sin u.\mathrm{d}u.\mathrm{d}\lambda.\mathrm{d}r
\]
Dabei muss die Integration nach $r$ von $r=0$ bis zu dem Wert an der Grenze von $t$ reichen, ferner von $\lambda=0$ bis $\lambda=2\pi$ und von $u=0$ bis $u=\pi$. Damit erhält $V$ notwendig einen bestimmten endlichen Wert.

\SourceLeaf{209}
Man sieht ferner leicht ein, dass man auch hier
\[
\frac{\mathrm{d}V}{\mathrm{d}x}=\int k\mathrm{d}t.\frac{\mathrm{d}\frac{1}{r}}{\mathrm{d}x}=\int\frac{k(a-x)\mathrm{d}t}{r^3}=X
\]
setzen darf. Das ist zulässig, weil sich auch dieser Ausdruck, der bei Verwendung von Polarkoordinaten in
\[
\int\!\int\!\int k\cos u.\sin u.\mathrm{d}u.\mathrm{d}\lambda.\mathrm{d}r
\]
übergeht, tatsächlich integrieren lässt. Also erhält $X$ einen bestimmten endlichen Wert, der sich stetig ändert, weil alle unendlich nahe bei $O$ liegenden Elemente nur einen unendlich kleinen Beitrag dazu liefern. Aus ähnlichen Gründen darf man auch
\[
\begin{array}{l}
\displaystyle\frac{\mathrm{d}V}{\mathrm{d}y}=\int\frac{k(b-y)\mathrm{d}t}{r^3}=Y\\[6pt]
\displaystyle\frac{\mathrm{d}V}{\mathrm{d}z}=\int\frac{k(c-z)\mathrm{d}t}{r^3}=Z
\end{array}
\]
setzen. Diese Größen haben also ebenso wie $V$ innerhalb von $t$ bestimmte Werte, die sich stetig ändern. Dasselbe gilt auch noch an der Grenze von $t$.

\Article{7}
Was nun aber die Ableitungen höherer Ordnung angeht, so muss man für Punkte innerhalb von $t$ anders vorgehen, da man zum Beispiel $\frac{\mathrm{d}X}{\mathrm{d}x}$ nicht in
\[
\int k\mathrm{d}t.\frac{\mathrm{d}\frac{a-x}{r^3}}{\mathrm{d}x}
\]
also in
\[
\int k(\frac{3(a-x)^2-rr}{r^5})\mathrm{d}t
\]
umformen darf. Bei genauer Betrachtung wäre dieser Ausdruck nur ein Zeichen ohne bestimmte, klare Bedeutung. Denn in jedem noch so kleinen Teil von $t$, der den Punkt einschließt, lassen sich Teilgebiete angeben, über die dieses Integral, wenn man es darauf erstreckt, jeden vorgegebenen Wert überschreitet, ob er positiv oder negativ ist. Damit fehlt hier die wesentliche
\SourceLeaf{210}
Bedingung, unter der allein das ganze Integral eine klare Bedeutung erhalten kann, nämlich die Anwendbarkeit der Exhaustionsmethode.

\Article{8}
Bevor wir diese Untersuchung in ihrer Allgemeinheit durchführen, ist es zur Klärung der Vorstellungen nützlich, einen sehr einfachen besonderen Fall zu betrachten.

Es sei $t$ eine Kugel mit dem Radius $=R$, deren Mittelpunkt mit dem Koordinatenursprung zusammenfällt. Die Dichte der Masse, die die Kugel ausfüllt, sei konstant $=k$, und den Abstand des Punktes $O$ vom Mittelpunkt bezeichnen wir mit $\rho=\BareRoot{(xx+yy+zz)}$. Bekanntlich hat das Potential zwei verschiedene Ausdrücke, je nachdem, ob $O$ innerhalb oder außerhalb der Kugel liegt. Im ersten Fall gilt
\[
V=2\pi kRR-\tfrac{2}{3}\pi k\rho\rho=2\pi kRR-\tfrac{2}{3}\pi k(xx+yy+zz)
\]
im zweiten dagegen
\[
V=\frac{4\pi kR^3}{3\rho}
\]
Auf der Kugeloberfläche ergeben beide Ausdrücke denselben Wert $\tfrac{4}{3}\pi kRR$, und das Potential ändert sich daher im ganzen Raum stetig.

Für die Ableitungen erhalten wir im Inneren
\[
\begin{array}{l}
\displaystyle\frac{\mathrm{d}V}{\mathrm{d}x}=X=-\tfrac{4}{3}\pi kx\\[6pt]
\displaystyle\frac{\mathrm{d}V}{\mathrm{d}y}=Y=-\tfrac{4}{3}\pi ky\\[6pt]
\displaystyle\frac{\mathrm{d}V}{\mathrm{d}z}=Z=-\tfrac{4}{3}\pi kz
\end{array}
\]
im Äußeren dagegen
\[
\begin{array}{l}
\displaystyle X=-\frac{4\pi kR^3x}{3\rho^3}\\[6pt]
\displaystyle Y=-\frac{4\pi kR^3y}{3\rho^3}\\[6pt]
\displaystyle Z=-\frac{4\pi kR^3z}{3\rho^3}
\end{array}
\]
Auch hier ergeben die letzteren Formeln auf der Oberfläche dieselben Werte wie die ersteren, sodass sich auch $X$, $Y$, $Z$ im ganzen Raum stetig ändern.

\SourceLeaf{211}
Mit den Ableitungen dieser Größen verhält es sich dagegen anders. Im Inneren haben wir
\[
\frac{\mathrm{d}X}{\mathrm{d}x}=-\tfrac{4}{3}\pi k,\qquad \frac{\mathrm{d}Y}{\mathrm{d}y}=-\tfrac{4}{3}\pi k,\qquad \frac{\mathrm{d}Z}{\mathrm{d}z}=-\tfrac{4}{3}\pi k
\]
im Äußeren dagegen
\[
\begin{array}{l}
\displaystyle\frac{\mathrm{d}X}{\mathrm{d}x}=\frac{4\pi kR^3(3xx-\rho\rho)}{3\rho^5}\\[6pt]
\displaystyle\frac{\mathrm{d}Y}{\mathrm{d}y}=\frac{4\pi kR^3(3yy-\rho\rho)}{3\rho^5}\\[6pt]
\displaystyle\frac{\mathrm{d}Z}{\mathrm{d}z}=\frac{4\pi kR^3(3zz-\rho\rho)}{3\rho^5}
\end{array}
\]
Auf der Oberfläche sind diese Werte jeweils um
\[
\frac{4\pi kxx}{RR},\qquad \frac{4\pi kyy}{RR},\qquad \frac{4\pi kzz}{RR}
\]
größer als die vorherigen. Die genannten Ableitungen ändern sich also im ganzen Inneren und im ganzen Äußeren stetig, aber beim Übergang von einem zum anderen sprunghaft. Auf der Trennfläche selbst muss man ihnen zwei Werte zuordnen, je nachdem, ob $\mathrm{d}x$, $\mathrm{d}y$, $\mathrm{d}z$ als positiv oder als negativ betrachtet werden.

Ähnliches gilt für die sechs übrigen Ableitungen
\[
\frac{\mathrm{d}X}{\mathrm{d}y},\qquad \frac{\mathrm{d}X}{\mathrm{d}z},\qquad \frac{\mathrm{d}Y}{\mathrm{d}x},\qquad \frac{\mathrm{d}Y}{\mathrm{d}z},\qquad \frac{\mathrm{d}Z}{\mathrm{d}x},\qquad \frac{\mathrm{d}Z}{\mathrm{d}y}
\]
die im Inneren der Kugel alle $=0$ sind und beim Durchgang durch die Kugeloberfläche sprunghaft die Änderungen
\[
\frac{4\pi kxy}{RR},\qquad \frac{4\pi kxz}{RR}\text{ u. s. f.}
\]
erfahren.

Die Summe $\frac{\mathrm{d}X}{\mathrm{d}x}+\frac{\mathrm{d}Y}{\mathrm{d}y}+\frac{\mathrm{d}Z}{\mathrm{d}z}$ oder $\frac{\mathrm{d}\mathrm{d}V}{\mathrm{d}x^2}+\frac{\mathrm{d}\mathrm{d}V}{\mathrm{d}y^2}+\frac{\mathrm{d}\mathrm{d}V}{\mathrm{d}z^2}$ wird im Inneren der Kugel $=-4\pi k$ und im Äußeren $=0$. Auf der Oberfläche selbst verliert sie jedoch ihre einfache Bedeutung. Genau genommen kann man nur sagen, dass sie eine Summe aus drei Teilen ist, von denen jeder zwei verschiedene Werte hat. Damit gibt es eigentlich acht Kombinationen. Eine davon stimmt mit dem Wert auf der Innenseite und eine andere mit dem Wert auf der Außenseite überein, während die übrigen sechs ohne jede Bedeutung bleiben. Die Ana%
\SourceLeaf{212}lyse, durch die einige Geometer auf der Kugeloberfläche den Wert $-2\pi k$ oder den Mittelwert zwischen dem inneren und dem äußeren Wert erhalten haben, kann ich nicht als zulässig anerkennen, sofern der Begriff des Differentialquotienten in seiner mathematischen Reinheit aufgefasst wird.

\Article{9}
Das Ergebnis des vorhergehenden Beispiels ist nur ein besonderer Fall des allgemeinen Satzes, nach dem, wenn der Punkt $O$ im Inneren der wirkenden Masse liegt, der Wert von $\frac{\mathrm{d}\mathrm{d}V}{\mathrm{d}x^2}+\frac{\mathrm{d}\mathrm{d}V}{\mathrm{d}y^2}+\frac{\mathrm{d}\mathrm{d}V}{\mathrm{d}z^2}$ gleich dem Produkt aus $-4\pi$ und der Dichte in $O$ wird. Die folgende Begründung dieses wichtigen Satzes scheint die befriedigendste zu sein.

Wir nehmen an, dass sich die Dichte $k$ innerhalb von $t$ nirgends sprunghaft ändert, oder dass sie eine mit $f(a,b,c)$ bezeichnete Funktion von $a$, $b$, $c$ ist, deren Wert sich innerhalb von $t$ überall stetig ändert und außerhalb von $t$ $=0$ wird.

Es sei $t'$ das Gebiet, in das $t$ übergeht, wenn die erste Koordinate jedes Punktes der Grenzfläche um $e$ vermindert wird, oder, was dasselbe ist, wenn man die Grenzfläche parallel zur ersten Koordinatenachse um $e$ zurückbewegt. Es bestehe $t$ aus den Gebieten $t^0$ und $\theta$ und $t'$ aus $t^0$ und $\theta'$, sodass $t^0$ das gesamte Gebiet ist, das $t$ und $t'$ gemeinsam bleibt. Wir betrachten die drei Integrale
\[
\begin{array}{ll}
\displaystyle\int\frac{f(a,b,c)(a-x)\mathrm{d}t}{((a-x)^2+(b-y)^2+(c-z)^2)^{\frac{3}{2}}}&.\ .\ .\ .\ .\ .\ .\ .\quad(1)\\[8pt]
\displaystyle\int\frac{f(a,b,c)(a-x-e)\mathrm{d}t}{((a-x-e)^2+(b-y)^2+(c-z)^2)^{\frac{3}{2}}}&.\ .\ .\ .\ .\ .\ .\quad(2)\\[8pt]
\displaystyle\int\frac{f(a+e,b,c)(a-x)\mathrm{d}t}{((a-x)^2+(b-y)^2+(c-z)^2)^{\frac{3}{2}}}&.\ .\ .\ .\ .\ .\ .\ .\quad(3)
\end{array}
\]
Dabei ist das Integral (1), über das ganze Gebiet $t$ erstreckt, der Wert von $\frac{\mathrm{d}V}{\mathrm{d}x}$ oder $X$ im Punkt $O$. Das Integral (2), ebenfalls über ganz $t$ erstreckt, ist der Wert von $\frac{\mathrm{d}V}{\mathrm{d}x}$ in dem Punkt mit den Koordinaten $x+e$, $y$, $z$. Diesen Wert wollen wir mit $X+\xi$ bezeichnen. Mit diesem Integral ist offenbar das Integral (3), über das ganze Gebiet $t'$ erstreckt, völlig identisch. Ist also
\SourceLeaf{213}%
\[
\begin{array}{rlcl}
\text{das Integral (1), erstreckt über}&t^0&.\ .\ .\ .\ .\ .\ .\ .&l\\
\text{über}&\theta&.\ .\ .\ .\ .\ .\ .\ .&\lambda\\
\text{das Integral (3), erstreckt über}&t^0&.\ .\ .\ .\ .\ .\ .\ .&l'\\
\text{über}&\theta'&.\ .\ .\ .\ .\ .\ .\ .&\lambda'
\end{array}
\]
so gilt $X=l+\lambda$, $X+\xi=l'+\lambda'$.

Setzen wir $f(a+e,b,c)-f(a,b,c)=\Delta k$, so ist das Integral
\[
\int\frac{\frac{\Delta k}{e}(a-x)\mathrm{d}t}{((a-x)^2+(b-y)^2+(c-z)^2)^{\frac{3}{2}}}\quad .\ .\ .\ .\ .\ .\ .\ .\quad(4)
\]
über $t^0$ erstreckt $=\frac{l'-l}{e}$.

Die bisherigen Ergebnisse gelten allgemein für jede Lage von $O$. Für die weitere Entwicklung schließen wir den Fall aus, dass $O$ auf der Oberfläche selbst liegt, oder nehmen an, dass $O$ in messbarer Entfernung von der Oberfläche innerhalb oder außerhalb von $t$ liegt.

Lassen wir nun $e$ unendlich klein werden, so sind die Gebiete $\theta$, $\theta'$ zwei unendlich schmale Raumschichten, die an der Oberfläche von $t$ anliegen. Zerlegen wir diese Oberfläche in Elemente $\mathrm{d}s$ und bezeichnen mit $\alpha$ den Winkel, den eine in $\mathrm{d}s$ nach außen errichtete Normale mit der ersten Koordinatenachse bildet, so ist $\alpha$ offenbar überall dort spitz, wo die Oberfläche von $t$ an $\theta$ grenzt, und dort stumpf, wo sie an $\theta'$ grenzt. Die Elemente von $\theta$ werden also durch $e\cos\alpha\mathrm{d}s$ ausgedrückt, die Elemente von $\theta'$ dagegen durch $-e\cos\alpha\mathrm{d}s$. Daraus lässt sich leicht schließen, dass $\frac{\lambda-\lambda'}{e}$ in das Integral
\[
\int\frac{f(a,b,c)(a-x)\cos\alpha.\mathrm{d}s}{((a-x)^2+(b-y)^2+(c-z)^2)^{\frac{3}{2}}}
\]
übergeht, oder, was dasselbe ist, in
\[
\int\frac{k(a-x)\cos\alpha.\mathrm{d}s}{r^3}
\]
über die ganze Oberfläche. Dabei bedeutet $k$ die Dichte am Element $\mathrm{d}s$.

Für einen unendlich kleinen Wert von $e$ geht ferner $\frac{\Delta k}{e}$ in den Wert der partiellen Ableitung $\frac{\mathrm{d}f(a,b,c)}{\mathrm{d}a}$ oder $\frac{\mathrm{d}k}{\mathrm{d}a}$ über und der Wert des Integrals (4) oder $\frac{(l'-l)}{e}$ in das Integral
\SourceLeaf{214}%
\[
\int\frac{\frac{\mathrm{d}k}{\mathrm{d}a}.(a-x)\mathrm{d}t}{r^3}
\]
über das ganze Gebiet $t$.

Schließlich ist für ein unendlich kleines $e$ der Ausdruck $\frac{l'-l}{e}-\frac{\lambda-\lambda'}{e}$ oder $\frac{\xi}{e}$ gerade der Wert der partiellen Ableitung $\frac{\mathrm{d}X}{\mathrm{d}x}$ oder $\frac{\mathrm{d}\mathrm{d}V}{\mathrm{d}x^2}$. Damit haben wir das einfache Ergebnis
\[
\frac{\mathrm{d}\mathrm{d}V}{\mathrm{d}x^2}=\frac{\mathrm{d}X}{\mathrm{d}x}=\int\frac{\frac{\mathrm{d}k}{\mathrm{d}a}.(a-x)\mathrm{d}t}{r^3}-\int\frac{k(a-x)\cos\alpha.\mathrm{d}s}{r^3}
\]
wobei die erste Integration über das ganze Gebiet $t$ und die zweite über dessen ganze Oberfläche zu erstrecken ist.

Dieses Ergebnis gilt, wie nahe $O$ auch auf der inneren oder äußeren Seite der Oberfläche liegen mag, sofern es nicht auf der Oberfläche selbst liegt. Dort hat $\frac{\mathrm{d}X}{\mathrm{d}x}$ vielmehr zwei verschiedene Werte. Das erste Integral ändert sich beim Durchgang durch die Oberfläche stetig. Dagegen ändert sich $-\int\frac{k(a-x)\cos\alpha\mathrm{d}s}{r^3}$ nach einem weiter unten zu beweisenden Satz beim Übergang von einem der Oberfläche unendlich nahen inneren Punkt zu einem äußeren um die endliche Größe $4\pi k\cos\alpha$. Dabei beziehen sich $k$ und $\alpha$ auf die Durchgangsstelle. Ebenso groß ist der Unterschied der beiden dort geltenden Werte von $\frac{\mathrm{d}X}{\mathrm{d}x}$.

\Article{10}
Auf ähnliche Weise erhalten wir, wenn $\beta$ und $\gamma$ in Bezug auf die zweite und dritte Koordinatenachse dieselbe Bedeutung haben wie $\alpha$ in Bezug auf die erste und für die Lage von $O$ dieselbe Einschränkung gilt wie zuvor,
\[
\begin{aligned}
\frac{\mathrm{d}Y}{\mathrm{d}y}&=\int\frac{\frac{\mathrm{d}k}{\mathrm{d}b}(b-y)\mathrm{d}t}{r^3}-\int\frac{k(b-y)\cos\beta.\mathrm{d}s}{r^3}\\
\frac{\mathrm{d}Z}{\mathrm{d}z}&=\int\frac{\frac{\mathrm{d}k}{\mathrm{d}c}(c-z)\mathrm{d}t}{r^3}-\int\frac{k(c-z)\cos\gamma.\mathrm{d}s}{r^3}
\end{aligned}
\]
Bedenken wir nun, dass
\[
\frac{\mathrm{d}k}{\mathrm{d}a}.\frac{a-x}{r}+\frac{\mathrm{d}k}{\mathrm{d}b}.\frac{b-y}{r}+\frac{\mathrm{d}k}{\mathrm{d}c}.\frac{c-z}{r}
\]
\SourceLeaf{215}%
gerade der Wert der Ableitung $\frac{\mathrm{d}k}{\mathrm{d}r}$ ist, sofern bei dieser Differentiation nur die Länge von $r$ als veränderlich und die Richtung als konstant betrachtet wird, und ferner, dass
\[
\frac{a-x}{r}.\cos\alpha+\frac{b-y}{r}.\cos\beta+\frac{c-z}{r}.\cos\gamma=\cos\psi
\]
gilt, wenn $\psi$ den Winkel zwischen der nach außen gerichteten Normalen in $\mathrm{d}s$ und der verlängerten Geraden $r$ bezeichnet, so folgt, dass, wenn das Integral
\[
\int\frac{\frac{\mathrm{d}k}{\mathrm{d}r}}{rr}.\mathrm{d}t
\]
über das ganze Gebiet $t$ mit $M$ und das Integral
\[
\int\frac{k\cos\psi}{rr}.\mathrm{d}s
\]
über die ganze Oberfläche von $t$ mit $N$ bezeichnet wird,
\[
\frac{\mathrm{d}\mathrm{d}V}{\mathrm{d}x^2}+\frac{\mathrm{d}\mathrm{d}V}{\mathrm{d}y^2}+\frac{\mathrm{d}\mathrm{d}V}{\mathrm{d}z^2}=M-N
\]
gilt.

Um die erste Integration auszuführen, beschreiben wir um den Mittelpunkt $O$ eine Kugelfläche mit dem Radius $1$ und zerlegen sie in Elemente $\mathrm{d}\sigma$. Die von $O$ durch alle Punkte des Randes von $\mathrm{d}\sigma$ gezogenen und unbegrenzt verlängerten Geraden bilden eine Kegelfläche (im weiteren Sinn des Wortes). Sie schneidet aus dem ganzen $t$ ein Gebiet heraus, das je nach den Umständen aus mehreren getrennten Stücken bestehen kann und dessen allgemeines Element $rr\mathrm{d}\sigma.\mathrm{d}r$ ist. Der Teil von $M$, der zu diesem Gebiet gehört, wird daher durch $\mathrm{d}\sigma.\int\frac{\mathrm{d}k}{\mathrm{d}r}.\mathrm{d}r$ ausgedrückt, wenn die Integration über alle in $t$ liegenden Teile einer durch $O$ und einen Punkt von $\mathrm{d}\sigma$ gehenden Geraden $r$ erstreckt wird, die so weit wie nötig verlängert ist. Nehmen wir nun an, diese Gerade schneide die Oberfläche von $t$ der Reihe nach in $O'$, $O''$, $O'''$, $O''''$ usw. Bezeichnen wir mit $r'$, $r''$, $r'''$, $r''''$ usw. die Werte von $r$ an diesen Punkten, mit $\mathrm{d}s'$, $\mathrm{d}s''$, $\mathrm{d}s'''$, $\mathrm{d}s''''$ usw. die entsprechenden Elemente, die der Elementarkegel aus der Oberfläche von $t$ herausschneidet, mit $k'$, $k''$, $k'''$, $k''''$ usw. die Werte von $k$ und mit $\psi'$, $\psi''$, $\psi'''$, $\psi''''$ usw. die Werte von $\psi$ an diesen Elementen, so sieht man leicht, dass
\SourceLeaf{216}%

I. im Fall, dass $O$ innerhalb von $t$ liegt, die Anzahl jener Punkte ungerade ist und die Integration $\int\frac{\mathrm{d}k}{\mathrm{d}r}.\mathrm{d}r$ von $r=0$ bis $r=r'$, dann von $r=r''$ bis $r=r'''$ usw. auszuführen ist. Bezeichnet man also die Dichte in $O$ mit $k^0$, so folgt
\[
\int\frac{\mathrm{d}k}{\mathrm{d}r}.\mathrm{d}r=-k^0+k'-k''+k'''-k''''+\text{ u. s. f.}
\]
Da die Winkel $\psi'$, $\psi''$, $\psi'''$, $\psi''''$ usw. offenbar abwechselnd spitz und stumpf sind, gilt
\[
\begin{aligned}
\mathrm{d}s'.\cos\psi'&=+r'r'\mathrm{d}\sigma\\
\mathrm{d}s''.\cos\psi''&=-r''r''\mathrm{d}\sigma\\
\mathrm{d}s'''.\cos\psi'''&=+r'''r'''\mathrm{d}\sigma\\
\mathrm{d}s''''.\cos\psi''''&=-r''''r''''\mathrm{d}\sigma
\end{aligned}
\]
usw. und damit
\[
\begin{aligned}
\mathrm{d}\sigma\int\frac{\mathrm{d}k}{\mathrm{d}r}.\mathrm{d}r&=-k^0\mathrm{d}\sigma+\frac{k'\cos\psi'}{r'r'}\mathrm{d}s'+\frac{k''\cos\psi''}{r''r''}\mathrm{d}s''+\frac{k'''\cos\psi'''}{r'''r'''}+\text{ u. s. f.}\\
&=-k^0\mathrm{d}\sigma+\sum\frac{k\cos\psi}{rr}\mathrm{d}s
\end{aligned}
\]
wobei über alle $\mathrm{d}s$ summiert wird, die dem Element $\mathrm{d}\sigma$ entsprechen. Durch Integration über sämtliche $\mathrm{d}\sigma$ erhält man also
\[
M=-4\pi k^0+\int\frac{k\cos\psi}{rr}\mathrm{d}s
\]
wobei das Integral über die ganze Oberfläche zu erstrecken ist, oder $M=-4\pi k^0+N$. Folglich gilt
\[
\frac{\mathrm{d}\mathrm{d}V}{\mathrm{d}x^2}+\frac{\mathrm{d}\mathrm{d}V}{\mathrm{d}y^2}+\frac{\mathrm{d}\mathrm{d}V}{\mathrm{d}z^2}=-4\pi k^0
\]

II. Im Fall, dass $O$ außerhalb von $t$ liegt, sind nur diejenigen $\mathrm{d}\sigma$ zu berücksichtigen, für die die Gerade durch $O$ und einen Punkt von $\mathrm{d}\sigma$ das Gebiet $t$ tatsächlich trifft. Die Anzahl der Punkte $O'$, $O''$, $O'''$ usw. ist hier stets gerade, und die Winkel $\psi'$, $\psi''$, $\psi'''$ usw. sind abwechselnd stumpf und spitz, also
\[
\begin{aligned}
\mathrm{d}s'.\cos\psi'&=-r'r'\mathrm{d}\sigma,\\
\mathrm{d}s''.\cos\psi''&=+r''r''\mathrm{d}\sigma,\\
\mathrm{d}s'''\cos\psi'''&=-r'''r'''\mathrm{d}\sigma\text{ u. s. f.}
\end{aligned}
\]
Da hier die Integration $\int\frac{\mathrm{d}k}{\mathrm{d}r}.\mathrm{d}r$ von $r=r'$ bis $r=r''$, dann von $r=r'''$
\SourceLeaf{217}%
bis $r=r''''$ usw. auszuführen ist, ergibt sich
\[
\begin{aligned}
\mathrm{d}\sigma.\int\frac{\mathrm{d}k}{\mathrm{d}r}.\mathrm{d}r&=\frac{k'\cos\psi'}{r'r'}.\mathrm{d}s'+\frac{k''\cos\psi''}{r''r''}.\mathrm{d}s''\\
&\quad{}+\frac{k'''\cos\psi'''}{r'''r'''}.\mathrm{d}s'''+\text{u. s. f.}\\
&=\sum\frac{k\cos\psi}{rr}\mathrm{d}s
\end{aligned}
\]
und nach der zweiten Integration über alle in Betracht kommenden $\mathrm{d}\sigma$
\[
M=\int\frac{k\cos\psi}{rr}\mathrm{d}s=N
\]
also, wie ohnehin bekannt ist,
\[
\frac{\mathrm{d}\mathrm{d}V}{\mathrm{d}x^2}+\frac{\mathrm{d}\mathrm{d}V}{\mathrm{d}y^2}+\frac{\mathrm{d}\mathrm{d}V}{\mathrm{d}z^2}=0
\]

\Article{11}
Obwohl unsere Beweisführung voraussetzt, dass sich die Dichte im \textit{ganzen} Gebiet $t$ stetig ändert, ist diese Bedingung für die Gültigkeit unseres Ergebnisses nicht notwendig. Es genügt, dass sich die Dichte im Punkt $O$ nach allen Richtungen stetig ändert oder dass $O$ innerhalb eines noch so kleinen Gebietes liegt, das diese Bedingung erfüllt. Setzen wir das Potential der in \textit{diesem} Gebiet enthaltenen Masse $=V'$ und das Potential der übrigen, außerhalb davon liegenden Massen $=V''$, so ist das gesamte Potential $V=V'+V''$. Da nach dem vorhergehenden Artikel
\[
\begin{aligned}
\frac{\mathrm{d}\mathrm{d}V'}{\mathrm{d}x^2}+\frac{\mathrm{d}\mathrm{d}V'}{\mathrm{d}y^2}+\frac{\mathrm{d}\mathrm{d}V'}{\mathrm{d}z^2}&=-4\pi k^0\\
\frac{\mathrm{d}\mathrm{d}V''}{\mathrm{d}x^2}+\frac{\mathrm{d}\mathrm{d}V''}{\mathrm{d}y^2}+\frac{\mathrm{d}\mathrm{d}V''}{\mathrm{d}z^2}&=0
\end{aligned}
\]
gilt, folgt
\[
\frac{\mathrm{d}\mathrm{d}V}{\mathrm{d}x^2}+\frac{\mathrm{d}\mathrm{d}V}{\mathrm{d}y^2}+\frac{\mathrm{d}\mathrm{d}V}{\mathrm{d}z^2}=-4\pi k^0
\]
Fehlt dagegen diese Bedingung im Punkt $O$ und liegt er also auf der Trennfläche zwischen zwei Gebieten, in denen sich die Dichte jeweils für sich stetig ändert, beim Übergang von einem zum anderen aber sprunghaft, so haben dort im Allgemeinen $\frac{\mathrm{d}\mathrm{d}V}{\mathrm{d}x^2}$, $\frac{\mathrm{d}\mathrm{d}V}{\mathrm{d}y^2}$, $\frac{\mathrm{d}\mathrm{d}V}{\mathrm{d}z^2}$ jeweils zwei verschiedene Werte. Für die Summe dieser Größen gilt dasselbe, was am Ende des 8. Artikels bemerkt wurde.
\SourceLeaf{218}%
\Article{12}
Wie bereits oben bemerkt, beziehen wir auch den idealen Fall in unsere Untersuchungen ein, bei dem die Anziehungs- oder Abstoßungskräfte von den Teilen einer \textit{Fläche} ausgehen sollen. Dabei erlauben wir uns die Darstellung, dass eine wirkende Masse auf der Fläche verteilt sei. Unter der Dichte in einem beliebigen Punkt der Fläche verstehen wir in diesem Fall den Quotienten aus der Masse, die in einem Flächenelement enthalten ist, dem dieser Punkt angehört, und diesem Element. Diese Dichte kann gleichförmig, also in allen Punkten gleich, oder ungleichförmig sein. Im letzteren Fall kann sie sich auf der ganzen Fläche stetig ändern, also in je zwei unendlich nahen Punkten nur unendlich wenig unterscheiden. Oder die Fläche kann in zwei oder mehr Stücke zerfallen, in denen jeweils eine stetige Änderung stattfindet, während sich die Dichte beim Übergang von einem Stück zum anderen sprunghaft ändert. Man kann sich auch eine Verteilung vorstellen, bei der die Dichte in einzelnen Punkten oder Linien unendlich groß wird, obwohl die gesamte Masse endlich bleibt. Der Fläche selbst wird, sofern sie keine Ebene ist, im Allgemeinen eine stetige Krümmung zugeschrieben, wobei Unterbrechungen in einzelnen Punkten (Ecken) oder Linien (Kanten) zugelassen bleiben.

\expandafter\def\csname HistoricalNoteText22\endcsname{Dass das Integral, das das Potential ausdrückt, endlich ist, erkennt man leicht, indem man die Fläche auf ähnliche Weise in Elemente zerlegt, wie es im 15. Artikel geschehen wird. Zugleich wird daraus ersichtlich, dass die Teile der Fläche, die den beiden fraglichen Punkten unendlich nahe liegen, zum ganzen Integral nur unendlich wenig beitragen. Daraus lässt sich das oben Gesagte leicht beweisen.}%
Unter diesen Voraussetzungen erhält das Potential auch in jedem Punkt der Fläche selbst, in dem die Dichte nicht unendlich groß ist, einen bestimmten endlichen Wert. Von diesem kann sich der Wert in einem zweiten Punkt, der ihm auf der Fläche oder außerhalb davon unendlich nahe liegt, nur unendlich wenig unterscheiden\HistoricalNote{22}. Mit anderen Worten ändert sich das Potential entlang jeder Linie stetig, ob sie auf der Fläche selbst liegt oder diese kreuzt.

\Article{13}
Bezeichnet man mit $k$ die Dichte im Flächenelement $\mathrm{d}s$, mit
\SourceLeaf{219}%
$a$, $b$, $c$ die Koordinaten eines zu ihm gehörenden Punktes, mit $r$ dessen Abstand von einem Punkt $O$ mit den Koordinaten $x$, $y$, $z$ und mit $V$ das Potential der in der Fläche enthaltenen Masse am Punkt $O$, so ist $V=\int\frac{k\mathrm{d}s}{r}$, wobei sich das Integral über die ganze Fläche erstreckt. Bezeichnet man schließlich mit $X$, $Y$, $Z$ die ebenso verstandenen Integrale
\[
\int\frac{k(a-x)\mathrm{d}s}{r^3},\quad\int\frac{k(b-y)\mathrm{d}s}{r^3},\quad\int\frac{k(c-z)\mathrm{d}s}{r^3}
\]
so sind $X$, $Y$, $Z$ zwar ganz gleichbedeutend mit $\frac{\mathrm{d}V}{\mathrm{d}x}$, $\frac{\mathrm{d}V}{\mathrm{d}y}$, $\frac{\mathrm{d}V}{\mathrm{d}z}$, solange $O$ außerhalb der Fläche liegt. Genau genommen gilt dies aber nicht mehr, wenn $O$ ein Punkt der Fläche selbst ist. Der Unterschied nimmt je nach dem Winkel, den die Flächennormale mit der betreffenden Koordinatenachse bildet, verschiedene Formen an. Es reicht offensichtlich aus, hier nur das Verhalten in Bezug auf die erste Koordinatenachse anzugeben.

I. Ist dieser Winkel $=0$, so hat in $O$ das Integral $X$ einen bestimmten Wert, $\frac{\mathrm{d}V}{\mathrm{d}x}$ dagegen zwei verschiedene Werte, je nachdem, ob man $\mathrm{d}x$ als positiv oder negativ betrachtet.

II. Ist der Winkel ein rechter, so lässt der Ausdruck für $X$ eine wirkliche Integration nicht zu (da dann eine ähnliche Bemerkung wie in Artikel 7 gilt), während $\frac{\mathrm{d}V}{\mathrm{d}x}$ nur einen bestimmten Wert hat.

III. Ist der Winkel spitz, so verhält sich $X$ wie im zweiten Fall und $\frac{\mathrm{d}V}{\mathrm{d}x}$ wie im ersten.

Weitere besondere Abwandlungen treten ein, wenn die Stetigkeit in $O$ entweder hinsichtlich der Dichte oder der Krümmung unterbrochen ist. Für unseren Hauptzweck ist es jedoch nicht nötig, solche Ausnahmefälle, die nur in einzelnen Linien oder Punkten auftreten können, ausführlich zu behandeln. Bei der näheren Untersuchung werden wir daher annehmen, dass im fraglichen Punkt eine bestimmte endliche Dichte und eine bestimmte Tangentialebene vorliegen.

\Article{14}
Bevor wir die Untersuchung in ihrer Allgemeinheit durchführen, wird es nützlich sein, einen einfachen besonderen Fall zu betrachten. Die Fläche sei ein Stück $A$ einer Kugelfläche, und die Dichte darin sei gleichförmig oder $k$ konstant. Es sind also $V$, $X$ die Werte der Integrale
\SourceLeaf{220}%
\[
\int\frac{k\mathrm{d}s}{r},\quad\int\frac{k(a-x)\mathrm{d}s}{r^3}
\]
über $A$. Bezeichnen wir mit $V'$, $X'$ dieselben Integrale über den übrigen Teil der Kugelfläche $B$ und mit $V^0$, $X^0$ die Integrale über die ganze Kugelfläche, so gilt $V=V^0-V'$, $X=X^0-X'$. Wir wollen außerdem den Radius der Kugel mit $R$ bezeichnen, den Koordinatenursprung in den Mittelpunkt der Kugel legen und $\BareRoot{(xx+yy+zz)}$, also den Abstand des Punktes $O$ vom Mittelpunkt der Kugel, $=\rho$ setzen.

Bekanntlich gilt $V^0=4\pi kR$, wenn $O$ innerhalb der Kugel liegt, dagegen $V^0=\frac{4\pi kRR}{\rho}$, wenn $O$ außerhalb liegt. Auf der Kugelfläche selbst fallen beide Werte zusammen. Der Differentialquotient $\frac{\mathrm{d}V^0}{\mathrm{d}x}$ ist daher innerhalb der Kugel $=0$ und außerhalb $=-\frac{4\pi kRRx}{\rho^3}$. Auf der Kugelfläche selbst gelten jedoch beide Werte zugleich, je nach dem Vorzeichen von $\mathrm{d}x$. Diese beiden Werte sind nur dann gleich, wenn $x=0$ ist, was dem Fall II des vorhergehenden Artikels entspricht.

Der Ausdruck für $X^0$, der innerhalb und außerhalb der Kugel mit $\frac{\mathrm{d}V^0}{\mathrm{d}x}$ gleichbedeutend ist, wird auf der Oberfläche zu einem leeren Zeichen, insofern eine wirkliche Integration unzulässig ist. Ausgenommen ist allein der Fall, in dem für die unendlich nahe liegenden Flächenelemente $a-x$ unendlich klein von höherer Ordnung als $r$ wird, nämlich wenn $y=0$, $z=0$. $x=\pm R$, für welchen Fall die Integration $X^0=\mp2\pi k$ ergibt, also keinen der Werte von $\frac{\mathrm{d}V^0}{\mathrm{d}x}$, sondern deren Mittelwert. Dieser Fall gehört offenbar zu I im vorhergehenden Artikel.

Bedenkt man nun, dass für einen Punkt $O$ auf der Kugeloberfläche innerhalb von $A$ die Größen $X'$ und $\frac{\mathrm{d}V'}{\mathrm{d}x}$ gleichbedeutend sind und bestimmte, sich stetig ändernde Werte haben, so wird deutlich, dass sich $X^0-X'$ und $\frac{\mathrm{d}V^0}{\mathrm{d}x}-\frac{\mathrm{d}V'}{\mathrm{d}x}$, also $X$ und $\frac{\mathrm{d}V}{\mathrm{d}x}$, zueinander genauso verhalten wie $X^0$ und $\frac{\mathrm{d}V^0}{\mathrm{d}x}$. Daraus folgen die im vorhergehenden Artikel aufgestellten Sätze von selbst.

\Article{15}
Für die allgemeinere Untersuchung ist es vorteilhaft, den Koordinatenursprung in einen Punkt $P$ auf der Fläche selbst zu legen und die erste Koordinatenachse senkrecht zur Tangentialebene in $P$ zu wählen.
\SourceLeaf{221}%

Bezeichnen wir mit $\psi$ den Winkel zwischen der Normalen auf ein beliebiges Flächenelement $\mathrm{d}s$ und der ersten Koordinatenachse, so ist $\cos\psi.\mathrm{d}s$ die Projektion von $\mathrm{d}s$ auf die Ebene von $b$ und $c$. Setzen wir
\[
\BareRoot{(bb+cc)}=\rho,\quad b=\rho\cos\theta,\quad c=\rho\sin\theta
\]
so stellt $\rho\mathrm{d}\rho.\mathrm{d}\theta$ ein beliebiges Element dieser Ebene dar, und das entsprechende Flächenelement ist $\mathrm{d}s=\frac{\rho\mathrm{d}\rho.\mathrm{d}\theta}{\cos\psi}$. Das darin enthaltene Massenelement ist also $=h\rho\mathrm{d}\rho.\mathrm{d}\theta$, wenn wir zur Abkürzung $h$ für $\frac{k}{\cos\psi}$ schreiben.

Wir wollen nun untersuchen, inwieweit sich der Wert von $X$ sprunghaft ändert, wenn der Punkt $O$ auf der ersten Koordinatenachse von einer Seite der Fläche auf die andere übergeht oder $x$ von einem negativen zu einem positiven Wert wechselt. Für diese Frage ist es offensichtlich gleichgültig, ob wir die ganze Fläche betrachten oder nur einen beliebig kleinen Teil, der den Punkt $P$ einschließt, da sich der Beitrag des übrigen Flächenteils zum Wert von $X$ stetig ändert. Wir dürfen daher $\rho$ nur von $0$ bis zu einem beliebig kleinen Grenzwert $\rho'$ laufen lassen und voraussetzen, dass sich innerhalb der so begrenzten Fläche $h$ und $\frac{a}{\rho}$ überall stetig ändern. Setzen wir für jeden bestimmten Wert von $\theta$ den Wert des Integrals $\int\frac{h(a-x)\rho\mathrm{d}\rho}{r^3}$ von $\rho=0$ bis $\rho=\rho'$ $=Q$, so wird $X=\int Q\mathrm{d}\theta$, wobei die Integration von $\theta=0$ bis $\theta=2\pi$ zu erstrecken ist.

Nun sind die Werte von $X$ für $x=0$, für ein unendlich kleines positives $x$ und für ein unendlich kleines negatives miteinander zu vergleichen (wobei die beiden anderen Koordinaten $y$, $z$ stets $=0$ angenommen werden). Wir bezeichnen diese drei Werte von $X$ mit $X^0$, $X'$, $X''$ und die entsprechenden Werte von $Q$ mit $Q^0$, $Q'$, $Q''$.

Da $r=\BareRoot{((a-x)^2+\rho\rho)}$, erhält man bei konstantem $\theta$
\[
\mathrm{d}\frac{h(a-x)}{r}=-\frac{h(a-x)\rho\mathrm{d}\rho}{r^3}+\frac{\mathrm{d}h}{\mathrm{d}\rho}.\frac{a-x}{r}.\mathrm{d}\rho+\frac{\mathrm{d}a}{\mathrm{d}\rho}.\frac{h\rho\rho}{r^3}.\mathrm{d}\rho
\]
und folglich
\[
Q=\int\frac{\mathrm{d}h}{\mathrm{d}\rho}.\frac{a-x}{r}.\mathrm{d}\rho+\int\frac{\mathrm{d}a}{\mathrm{d}\rho}.\frac{h\rho\rho}{r^3}.\mathrm{d}\rho-\frac{h'(a'-x)}{r'}+\text{Const.}
\]
wobei beide Integrationen von $\rho=0$ bis $\rho=\rho'$ reichen und die Werte
\SourceLeaf{222}%
von $h,a,r$ für $\rho=\rho'$ mit $h',a',r'$ bezeichnet sind. Als Konstante ist der Wert von $\frac{h(a-x)}{r}$ für $\rho=0$ zu nehmen. Bezeichnet man die Dichte in $P$ mit $k^0$, so ist dieser Wert $=-k^0$ für ein positives $x$ und $=+k^0$ für ein negatives, da für $\rho=0$ offensichtlich $a=0$, $\psi=0$, $h=k^0$, $x=\pm r$ gilt. Im Fall $x=0$ ist dagegen als Konstante der Grenzwert von $\frac{ha}{r}$ bei unendlich abnehmendem $\rho$ zu nehmen. Er ist $=0$, weil $a$ unendlich klein von höherer Ordnung als $r$ wird.

Der Wert des Integrals $\int\frac{\mathrm dh}{\mathrm d\rho}.\frac{a-x}{r}.\mathrm d\rho$ bleibt bis auf einen unendlich kleinen Unterschied derselbe, ob man $x=0$ oder unendlich klein $=\pm\varepsilon$ setzt. Zerlegt man nämlich dieses Integral in
\[
\int_0^\delta\frac{\mathrm dh}{\mathrm d\rho}.\frac{a-x}{r}.\mathrm d\rho+\int_\delta^{\rho'}\frac{\mathrm dh}{\mathrm d\rho}.\frac{a-x}{r}.\mathrm d\rho
\]
so gilt die Behauptung offensichtlich für den ersten Teil, wenn $\delta$ unendlich klein ist, und für den zweiten, wenn $\frac{\delta}{\varepsilon}$ unendlich groß ist, also für das Ganze, wenn $\delta$ unendlich klein von niedrigerer Ordnung als $\varepsilon$ ist.

Ein ähnlicher Schluss gilt auch für das Integral $\int\frac{\mathrm da}{\mathrm d\rho}.\frac{h\rho\rho}{r^3}.\mathrm d\rho$, wenn die Punkte der Fläche, die einem bestimmten Wert von $\theta$ entsprechen, eine Kurve mit messbarer Krümmung in $P$ bilden, sodass $\frac{a}{\rho\rho}$ im hier betrachteten Gebiet einen endlichen, sich stetig ändernden Wert hat. Bezeichnet man diesen Wert mit $A$, so wird
\[
\frac{\mathrm da}{\mathrm d\rho}=2A\rho+\frac{\mathrm dA}{\mathrm d\rho}.\rho\rho
\]
und das Integral zerfällt in die folgenden zwei
\[
\int\frac{2\rho^3Ah\mathrm d\rho}{r^3}+\int\frac{\mathrm dA}{\mathrm d\rho}.\frac{\rho^4}{r^3}h\mathrm d\rho
\]
für die beide die Gültigkeit der obigen Schlussweise unmittelbar einleuchtet.

Schließlich sind auch die Werte von $\frac{h'(a'-x)}{r'}$ für alle drei Werte von $x$ bis auf unendlich kleine Unterschiede offensichtlich gleich.

Daraus folgt also, dass $Q'+k^0$, $Q^0$, $Q''-k^0$ bis auf unendlich kleine Unterschiede gleich sind. Dasselbe gilt demnach auch für
\[
\int(Q'+k^0)\mathrm d\theta,\quad\int Q^0\mathrm d\theta,\quad\int(Q''-k^0)\mathrm d\theta
\]
oder für die Größen
\[
X'+2\pi k^0,\quad X^0,\quad X''-2\pi k^0
\]

\SourceLeaf{223}%
Diesen wichtigen Satz kann man auch so ausdrücken. Der Grenzwert von $X$ bei unendlich abnehmendem positivem $x$ ist $X^0-2\pi k^0$, bei unendlich abnehmendem negativem $x$ dagegen $X^0+2\pi k^0$. Mit anderen Worten, $X$ ändert sich zweimal sprunghaft um $-2\pi k^0$, wenn $x$ von einem negativen zu einem positiven Wert übergeht, das erste Mal, wenn $x$ den Wert $0$ erreicht, und das zweite Mal, wenn es ihn überschreitet.

\Article{16}
Im Beweis des vorhergehenden Artikels wird vorausgesetzt, dass die Schnitte der Fläche mit den durch die erste Koordinatenachse gelegten Ebenen in $P$ eine messbare Krümmung haben. Unser Ergebnis gilt aber mit einer einzigen Ausnahme auch dann noch, wenn die Krümmung in $P$ unendlich groß ist. Dass $\frac a\rho$ für unendlich kleines $\rho$ selbst unendlich klein werden muss, folgt schon aus der Voraussetzung, dass die Fläche in $P$ eine bestimmte Tangentialebene hat. Die beiden Größen sind jedoch nur dann von gleicher Ordnung, wenn der Krümmungsradius endlich ist. Bei einem unendlich kleinen Krümmungsradius ist $\frac a\rho$ dagegen von niedrigerer Ordnung als $\rho$. Wir werden nun zeigen, dass unsere Ergebnisse auch im letzteren Fall gültig bleiben, sofern die Ordnungen beider Größen \textit{vergleichbar} sind.

Nehmen wir also an, dass $\frac a\rho$ von derselben Ordnung ist wie $\rho^\mu$, wobei $\mu$ einen endlichen positiven Exponenten bezeichnet. Dann ist $\frac{a}{\rho^{1+\mu}}$ eine endliche Größe, die sich im betrachteten Gebiet stetig ändert und die wir mit $B$ bezeichnen wollen. Das Integral $\int\frac{\mathrm da}{\mathrm d\rho}.\frac{h\rho\rho}{r^3}\mathrm d\rho$ zerfällt also in die folgenden beiden Integrale
\[
\int\frac{(1+\mu)\rho^{2+\mu}hB\mathrm d\rho}{r^3}+\int\frac{\rho^{3+\mu}}{r^3}.\frac{\mathrm dB}{\mathrm d\rho}.h\mathrm d\rho
\]
Auf das zweite Integral lassen sich die Schlüsse des vorhergehenden Artikels unmittelbar anwenden, auf das erste nach einer kleinen Umformung. Setzt man nämlich $\frac1\mu=m$, $\rho^\mu=\sigma$ oder $\rho=\sigma^m$, so wird dieses Integral
\[
=(m+1)\int\frac{Bh\sigma^{3m}\mathrm d\sigma}{(\sigma^{2m}+(a-x)^2)^{\frac32}}
\]
Auch dieses Integral hat offensichtlich nur einen unendlich kleinen Wert, solange die Integration nur von $0$ bis zu einem unendlich kleinen Wert von $\sigma$ reicht. Für jeden endlichen Wert von $\sigma$ hingegen hat der Koeffizient
\SourceLeaf{224}%
von $\mathrm d\sigma$ bis auf einen unendlich kleinen Unterschied denselben Wert, ob man $x=0$ oder unendlich klein annimmt. Das gilt also auch für das ganze Integral, wenn man von $\sigma=0$ bis $\sigma=\BareRoot[m]{\rho'}$ integriert.

Unsere Schlüsse verlieren nur in einem einzigen Fall ihre Gültigkeit, nämlich wenn $\frac a\rho$ mit keiner Potenz von $\rho$ mehr von gleicher Ordnung ist, etwa wenn $\frac a\rho$ von derselben Ordnung wäre wie $\frac{1}{\log\frac1\rho}$. In diesem Fall würde $Q$ bei unendlicher Annäherung des Punktes $O$ an die Fläche über jede Grenze wachsen. Dasselbe würde auch für $X$ gelten, wenn dieses Verhalten bei allen Werten von $\theta$ aufträte und nicht nur bei einem oder einigen. Wir brauchen das hier jedoch nicht weiter auszuführen, da wir diesen singulären Fall ohne Nachteil ganz aus unserer Untersuchung ausschließen können.

\Article{17}
Wir betrachten nun unter denselben Voraussetzungen und mit denselben Bezeichnungen wie in Artikel 15 die Größe $Y$, von der $\frac{hb\mathrm db.\mathrm dc}{r^3}$ ein beliebiges Element ist. Da $r=\BareRoot{(bb+cc+(a-x)^2)}$, gilt
\[
\frac{\mathrm d\frac hr}{\mathrm db}=-\frac{hb}{r^3}+\frac1r.\frac{\mathrm dh}{\mathrm db}-\frac{h(a-x)}{r^3}.\frac{\mathrm da}{\mathrm db}
\]
wenn $c$ konstant gehalten wird. Die erste Integration unter dieser Voraussetzung ergibt daher
\[
\int\frac{hb\mathrm db}{r^3}=\frac{h^*}{r^*}-\frac{h^{**}}{r^{**}}+\int\frac1r.\frac{\mathrm dh}{\mathrm db}.\mathrm db-\int\frac{h(a-x)}{r^3}.\frac{\mathrm da}{\mathrm db}.\mathrm db
\]
wobei die Integrationen vom kleinsten bis zum größten Wert von $b$ für jeden bestimmten Wert von $c$ reichen und mit $h^*,r^*,h^{**},r^{**}$ die zu diesen Grenzen gehörenden Werte von $h$ und $r$ bezeichnet sind, Schreiben wir zur Abkürzung
\[
\frac{h^*}{r^*}-\frac{h^{**}}{r^{**}}=T,\quad\frac\rho r.\frac{\mathrm dh}{\mathrm db}-\frac{h(a-x)\rho}{r^3}.\frac{\mathrm da}{\mathrm db}=U
\]
so ergibt sich
\[
Y=\int T\mathrm dc+\iint\frac U\rho.\mathrm db.\mathrm dc
\]
\SourceLeaf{225}%
wobei die Integration nach $c$ vom kleinsten Wert dieser Koordinate auf der Fläche bis zum größten reichen muss. Im Doppelintegral stellt $\mathrm db.\mathrm dc$ die Projektion eines beliebigen Flächenelements auf die Ebene der Koordinaten $b,c$ dar. Man kann dafür also auch $\rho\mathrm d\rho.\mathrm d\theta$ schreiben. Damit ergibt sich
\[
Y=\int T\mathrm dc+\iint U\mathrm d\rho.\mathrm d\theta
\]
wobei im Doppelintegral von $\rho=0$ bis $\rho=\rho'$ und von $\theta=0$ bis $\theta=2\pi$ zu integrieren ist. Mit ähnlichen Überlegungen wie in Artikel 15 erkennt man leicht, dass dieser Ausdruck bis auf unendlich kleine Unterschiede denselben Wert annimmt, ob man $x=0$ oder unendlich klein ansetzt. Mit anderen Worten hat der Wert von $Y$ bei unendlich klein werdenden positiven und negativen Werten von $x$ denselben Grenzwert. Das ist genau der Wert der obigen Formel, wenn man darin $x=0$ setzt. Entsprechend wollen wir diesen Wert mit $Y^0$ bezeichnen. Dabei ist jedoch zu beachten, dass man ihn nicht \textit{den} Wert von $\int\frac{kb\mathrm ds}{r^3}$ für $x=0$ nennen darf (da dieser Ausdruck für $x=0$ keine wirkliche Integration zulässt). Es ist nur \textit{ein} Wert dieses Integrals, nämlich derjenige, der sich bei der oben verwendeten Reihenfolge der Integration ergibt.

Dieses Ergebnis bedarf ähnlich wie in Artikel 16 einer Einschränkung im singulären Fall, dass die Krümmungsradien im Punkt $P$ unendlich klein sind, ebenso wenn dort $\frac{\mathrm dh}{\mathrm db}$ unendlich groß wird. Für unseren Zweck ist es jedoch nicht nötig, diese Ausnahmefälle eigens zu betrachten. Sie können nur in einzelnen Punkten oder Linien auftreten (also nicht innerhalb von Teilen der Fläche, sondern nur an deren Grenzen).

Schließlich ist unmittelbar klar, dass sich die Größe $Z$ oder das Integral $\int\frac{kc\mathrm ds}{r^3}$ genauso verhält wie $Y$. Nähert sich der Punkt $O$ auf der ersten Koordinatenachse dem Punkt $P$ bis auf einen unendlich kleinen Abstand, so hat dieses Integral denselben Grenzwert $Z^0$, ob die Annäherung von der positiven oder der negativen Seite erfolgt. Dieser Grenzwert ist zugleich der Wert von $\iint\frac{hc\mathrm dc.\mathrm db}{r^3}$ für $x=0$, wenn man zuerst nach $c$ integriert.

\SourceLeaf{226}%


\Article{18}
Bedenken wir nun, dass die Größen $\frac{\mathrm dV}{\mathrm dx}$, $\frac{\mathrm dV}{\mathrm dy}$, $\frac{\mathrm dV}{\mathrm dz}$ an allen Punkten des Raumes außerhalb der Fläche selbst uneingeschränkt mit $X,Y,Z$ übereinstimmen und dass sich $V$ überall stetig ändert. Dann folgt aus den Ergebnissen des vorhergehenden Artikels leicht, dass sich in unendlich kleinem Abstand von $P$, also für unendlich kleine Werte von $x,y,z$, der Wert von $V$ bis auf unendlich kleine Größen höherer Ordnung ausdrücken lässt durch
\[
V^0+x(X^0-2\pi k^0)+yY^0+zZ^0
\]
wenn $x$ positiv ist, oder durch
\[
V^0+x(X^0+2\pi k^0)+yY^0+zZ^0
\]
wenn $x$ negativ ist. Dabei bezeichnet $V^0$ den Wert von $V$ im Punkt $P$ selbst, also für $x=0$, $y=0$, $z=0$. Betrachten wir also die Werte von $V$ auf einer Geraden durch $P$, die mit den drei Koordinatenachsen die Winkel $A,B,C$ bildet. Wir bezeichnen mit $t$ ein beliebiges Stück dieser Geraden und mit $t^0$ den Wert von $t$ im Punkt $P$. Ist $t-t^0$ unendlich klein, so gilt bis auf eine unendlich kleine Größe höherer Ordnung
\[
V=V^0+(t-t^0)(X^0\cos A+Y^0\cos B+Z^0\cos C\mp2\pi k^0\cos A)
\]
wobei das obere Zeichen für positive und das untere für negative Werte von $(t-t^0)\cos A$ gilt. Anders gesagt hat $\frac{\mathrm dV}{\mathrm dt}$ im Punkt $P$ für einen spitzen Winkel $A$ zwei verschiedene Werte, nämlich
\[
\begin{gathered}
X^0\cos A+Y^0\cos B+Z^0\cos C-2\pi k^0\cos A\text{ und}\\
X^0\cos A+Y^0\cos B+Z^0\cos C+2\pi k^0\cos A
\end{gathered}
\]
je nachdem, ob man $\mathrm dt$ als positiv oder negativ ansieht. Ist $A$ ein rechter Winkel, sodass die Gerade die Fläche nur berührt, so fallen beide Ausdrücke zusammen, und es ergibt sich
\[
\frac{\mathrm dV}{\mathrm dt}=Y^0\cos B+Z^0\cos C
\]
\begin{center}*\qquad\qquad*\\{}*\end{center}

\SourceLeaf{227}%
Die bisher vorgetragenen Sätze sind ihrem wesentlichen Inhalt nach zwar nicht neu, durften aber um des Zusammenhangs willen. als notwendige Vorbereitungen für die folgenden Untersuchungen nicht übergangen werden, in. denen eine Reihe neuer Lehrsätze entwickelt werden wird.

\Article{19}
Es sei $V$ das Potential eines Systems von Massen $M',M'',M'''\ldots$ in den Punkten $P',P'',P'''\ldots$. Ferner sei $v$ das Potential eines zweiten Systems von Massen $m',m'',m'''\ldots$ in den Punkten $p',p'',p'''\ldots$. Es seien weiter $V',V'',V'''\ldots$ die Werte von $V$ in den letzteren Punkten und $v',v'',v'''\ldots$ die Werte von $v$ in den Punkten $P',P'',P'''\ldots$ Dann gilt die Gleichung
\[
M'v'+M''v''+M'''v'''+\text{u. s. f.}=m'V'+m''V''+m'''V'''+\text{u. s. f.}
\]
die sich auch als $\sum Mv=\sum mV$ schreiben lässt, wobei $M$ eine beliebige Masse des ersten und $m$ eine beliebige Masse des zweiten Systems bezeichnet. Tatsächlich ist sowohl $\sum Mv$ als auch $\sum mV$ die Summe aller Kombinationen $\frac{Mm}{\rho}$, wenn $\rho$ den Abstand zwischen den Punkten bezeichnet, in denen sich die betreffenden Massen $M,m$ befinden.

Sind die Massen des einen Systems oder beider Systeme stetig auf Linien, auf Flächen oder in räumlichen Gebieten verteilt und nicht an einzelnen Punkten konzentriert, so bleibt die obige Gleichung gültig, wenn man an die Stelle der Summe das entsprechende Integral setzt.

Ist also zum Beispiel das zweite Massensystem so auf einer Fläche verteilt, dass auf das Flächenelement $\mathrm ds$ die Masse $k\mathrm ds$ entfällt, so gilt $\sum Mv=\int kV\mathrm ds$. Gilt Entsprechendes auch für das erste System, sodass das Flächenelement $\mathrm dS$ die Masse $K\mathrm dS$ enthält, so ergibt sich $\int Kv\mathrm dS=\int kV\mathrm ds$. Für den letzten Fall ist es wichtig zu beachten, dass diese Gleichung auch dann noch gilt, wenn beide Flächen zusammenfallen. Der Kürze wegen wollen wir hier nur die wesentlichen Schritte andeuten, mit denen sich diese Erweiterung des Satzes streng begründen lässt. Es ist nämlich nicht schwer zu zeigen, dass diese beiden Integrale, soweit sie sich auf ein und dieselbe Fläche beziehen, die Grenzwerte der Integrale für zwei getrennte Flächen
\SourceLeaf{228}%
sind, wenn man deren Abstand unendlich klein werden lässt. Dazu braucht man die beiden Flächen nur gleich und parallel anzunehmen. Unmittelbar einleuchtend ist dieses Beweisverfahren allerdings nur, wenn die gegebene Fläche so beschaffen ist, dass die Normalen in allen ihren Punkten mit ein und derselben Geraden spitze Winkel bilden. Bei einer Fläche, die diese Bedingung nicht erfüllt (wie stets bei einer geschlossenen Fläche), muss man die Fläche zunächst in zwei oder mehrere Teile zerlegen, die diese Bedingung jeweils erfüllen. So lässt sich dieser Fall leicht auf den vorherigen zurückführen.

\Article{20}
Wenden wir den Satz des vorhergehenden Artikels auf den Fall an, dass das zweite Massensystem mit gleichmäßiger Dichte $k=1$ auf einer Kugelfläche mit Radius $=R$ verteilt ist. Dann ist das daraus entstehende Potential $v$ im Inneren der Kugel konstant $=4\pi R$. An jedem Punkt außerhalb der Kugel mit Abstand $=r$ vom Mittelpunkt gilt $v=\frac{4\pi RR}{r}$. Das ist ebenso groß wie das Potential im Mittelpunkt, das von einer in jenem Punkt angenommenen Masse $4\pi RR$ herrührt. Auf der Kugeloberfläche fallen beide Werte von $v$ zusammen. Liegt also das erste Massensystem ganz im Inneren der Kugel, so ist $\sum Mv$ äqual dem Produkt aus der Gesamtmasse dieses Systems und $4\pi R$. Liegt jenes Massensystem ganz außerhalb der Kugel, so ist $\sum Mv$ äqual dem Produkt aus dem Potential dieser Masse im Kugelmittelpunkt und $4\pi RR$. Ist schließlich das erste Massensystem stetig auf der Kugeloberfläche verteilt, so kann man für $\int Kv\mathrm dS$ beide Ausdrücke gleichermaßen verwenden. Daraus folgt der folgende Lehrsatz.

\textsc{Lehrsatz.} Bezeichnet $V$ das Potential einer beliebig verteilten Masse in einem Element einer Kugelfläche mit Radius $R$, das wir mit $\mathrm ds$ bezeichnen, so ergibt die Integration über die ganze Kugelfläche
\[
\int V\mathrm ds=4\pi(RM^0+RRV^0)
\]
wenn $M^0$ die gesamte im Inneren der Kugel befindliche Masse und $V^0$ das Potential der außerhalb liegenden Masse im Kugelmittelpunkt bezeichnet. Dabei kann man die Massen, die etwa stetig auf der Kugeloberfläche verteilt sind, nach Belieben den äußeren oder den inneren Massen zuordnen.

\SourceLeaf{229}%
\Article{21}
\textsc{Lehrsatz.} Das Potential $V$ von Massen, die sämtlich außerhalb eines zusammenhängenden Gebiets liegen, kann nicht in einem Teil dieses Gebiets einen konstanten Wert und zugleich in einem anderen Teil desselben Gebiets einen anderen Wert haben.

\textit{Beweis.} Nehmen wir an, das Potential sei in jedem Punkt des Gebiets $A$ konstant $=a$ und in jedem Punkt eines anderen, an $A$ grenzenden und keine Masse enthaltenden Gebiets $B$ (algebraisch) größer als $a$. Man konstruiere eine Kugel, die zu einem Teil in $B$ und mit ihrem übrigen Teil samt dem Mittelpunkt in $A$ liegt. Eine solche Konstruktion ist stets möglich. Ist nun $R$ der Radius dieser Kugel und $\mathrm ds$ ein beliebiges Element ihrer Oberfläche, so gilt nach dem Lehrsatz des vorhergehenden Artikels $\int V\mathrm ds=4\pi RRa$ und $\int(V-a)\mathrm ds=0$. Das ist unmöglich, denn für den in $A$ liegenden Teil der Oberfläche gilt $V-a=0$, und für den übrigen Teil ist die Differenz nach Voraussetzung nicht $=0$, sondern positiv.

Ganz entsprechend sieht man, dass es unmöglich ist, dass in allen Punkten eines an $A$ grenzenden Gebiets $V$ kleiner als $a$ ist.

Mindestens einer dieser beiden Fälle müsste aber offenbar eintreten, wenn unser Lehrsatz falsch wäre.

Dieser Lehrsatz enthält die folgenden beiden Aussagen.

I. Umschließt das Gebiet, in dem sich die Massen befinden, schalenförmig ein massenleeres Gebiet und hat das Potential in einem Teil dieses eingeschlossenen Gebiets einen konstanten Wert, so gilt dieser Wert für alle Punkte des gesamten eingeschlossenen Gebiets.

II. Hat das Potential der in einem endlichen Gebiet eingeschlossenen Massen in irgendeinem Teil des äußeren Gebiets einen konstanten Wert, so gilt dieser Wert für das gesamte unendliche äußere Gebiet.

Zugleich sieht man leicht, dass der konstante Wert des Potentials in diesem zweiten Fall nur $0$ sein kann. Bezeichnet man nämlich mit $M$ die Summe aller Massen, wenn sie sämtlich dasselbe Vorzeichen haben, oder andernfalls nur die Summe der positiven oder der negativen Massen, je nachdem, welche von beiden überwiegen, so ist das Potential in einem Punkt mit Abstand $=r$ vom nächsten Massenelement jedenfalls, abso%
\SourceLeaf{230}%
lut genommen, kleiner als $\frac Mr$. Dieser Bruch kann im äußeren Gebiet offenbar kleiner als jede vorgegebene Größe werden.

\Article{22}
\textsc{Lehrsatz.} Es sei $\mathrm ds$ das Element einer Fläche, die ein zusammenhängendes endliches Gebiet begrenzt, und $P$ die Kraft, die beliebig verteilte Massen in $\mathrm ds$ in der zur Fläche normalen Richtung ausüben. Dabei gilt eine nach innen oder außen gerichtete Kraft als positiv, je nachdem, ob anziehende oder abstoßende Massen als positiv gelten. Dann ist das über die ganze Fläche erstreckte Integral $\int P\mathrm ds$ $=4\pi M+2\pi M'$, wobei $M$ die Summe der im Inneren des Gebiets befindlichen Massen bezeichnet und $M'$ die Summe der stetig auf der Oberfläche verteilten Massen.

\textit{Beweis.} Bezeichnet man mit $U\mathrm d\mu$ den Anteil von $P$, der vom Massenelement $\mathrm d\mu$ herrührt, mit $r$ den Abstand des Elements $\mathrm d\mu$ von $\mathrm ds$ und mit $u$ den Winkel, den die nach innen gerichtete Normale in $\mathrm ds$ mit $r$ bildet, so gilt $U=\frac{\cos u}{rr}$. Für jedes bestimmte $\mathrm d\mu$ gilt aber nach einem in Artikel 6 der \textit{Theoria Attractionis corporum sphaeroidicorum ellipticorum} bewiesenen Lehrsatz $\int\frac{\cos u}{rr}.\mathrm ds=0$, $2\pi$ oder $4\pi$, je nachdem, ob $\mathrm d\mu$ außerhalb des von der Fläche begrenzten Gebiets, in der Fläche selbst oder innerhalb jenes Gebiets liegt. Da nun $\int P\mathrm ds$ die Summe aller Beiträge $\mathrm d\mu.\int U\mathrm ds$ ist, folgt unser Lehrsatz unmittelbar.

Zu dem hier verwendeten Hilfssatz ist noch zu bemerken, dass er in der Form, in der er an der angeführten Stelle ausgesprochen ist, für einen besonderen Fall einer Abänderung bedarf. Es bezeichnet nämlich $r$ den Abstand eines \textit{gegebenen Punktes} vom Element $\mathrm ds$. Liegt dieser Punkt in der Fläche selbst, so gilt die Formel $\int\frac{\cos u}{rr}.\mathrm ds=2\pi$ nur, wenn die Krümmung der Fläche in diesem Punkt stetig ist. Diese Stetigkeit ist jedoch unterbrochen, wenn der Punkt auf einer Kante oder an einer Ecke liegt. Dann ist anstelle von $2\pi$ der Flächeninhalt der Figur einzusetzen, die alle von diesem Punkt ausgehenden und die Fläche berührenden Geraden aus einer Kugeloberfläche mit Radius $1$ um diesen Punkt als Mittelpunkt ausschneiden. Da solche Ausnahmefälle jedoch nur Linien oder Punkte betreffen, also keine \textit{Teile} der Fläche, sondern nur Scheidungs%
\SourceLeaf{231}%
grenzen zwischen Teilen, hat dies auf die hier gemachte Anwendung des Hilfssatzes offenbar keinerlei Einfluss.

\Article{23}
Wir legen durch jeden Punkt der Fläche eine Normale und bezeichnen mit $p$ den Abstand eines beliebigen Punktes auf dieser Normalen von ihrem in der Fläche selbst liegenden Ursprung. Auf der inneren Seite der Fläche gilt der Abstand als positiv. Das Potential $V$ der Massen lässt sich als Funktion von $p$ und zwei weiteren veränderlichen Größen auffassen, die die einzelnen Punkte der Fläche auf irgendeine Weise voneinander unterscheiden. Ebenso verhält es sich mit der partiellen Ableitung $\frac{\mathrm dV}{\mathrm dp}$, deren Wert hier aber nur für Punkte in der Fläche selbst, also für $p=0$, betrachtet werden soll. Diese Ableitung ist mit $P$ völlig gleichbedeutend, wenn sich Massen nur im inneren Gebiet, nur im äußeren oder in beiden befinden und keine Masse auf der Fläche selbst verteilt ist. Dann gilt
\[
\int\frac{\mathrm dV}{\mathrm dp}.\mathrm ds=4\pi M
\]
Ist dagegen die ganze Masse allein auf der Fläche selbst verteilt, sodass das Element $\mathrm ds$ die Masse $k\mathrm ds$ enthält, so sind $\frac{\mathrm dV}{\mathrm dp}$ und $P$ nicht mehr gleichbedeutend. Die letztere Größe hat hier in Bezug auf $p$ offenbar dieselbe Bedeutung wie $X^0$ in Bezug auf $x$ in Artikel 15. Die Ableitung $\frac{\mathrm dV}{\mathrm dp}$ hat dagegen zwei verschiedene Werte, nämlich $P-2\pi k$ und $P+2\pi k$, je nachdem, ob $\mathrm dp$ als positiv oder negativ betrachtet wird. Nun ist $\int k\mathrm ds$ offensichtlich gleich der gesamten auf der Fläche verteilten Masse $M'$, und nach dem Lehrsatz des vorhergehenden Artikels gilt $\int P\mathrm ds=2\pi M'$. Damit erhält man
\[
\int\frac{\mathrm dV}{\mathrm dp}.\mathrm ds=0\quad\text{oder}\quad\int\frac{\mathrm dV}{\mathrm dp}.\mathrm ds=4\pi M'
\]
je nachdem, ob man für $\frac{\mathrm dV}{\mathrm dp}$ überall den auf der inneren oder den auf der äußeren Seite der Fläche geltenden Wert versteht. Das Integral $\int\frac{\mathrm dV}{\mathrm dp}.\mathrm ds$ verhält sich also im ersten Fall genauso, als gehörte die Masse $M'$ zum äußeren Gebiet, und im zweiten Fall, als gehörte sie zum inneren.

Bei beliebig verteilten Massen gilt daher die Gleichung
\[
\int\frac{\mathrm dV}{\mathrm dp}.\mathrm ds=4\pi M
\]
\SourceLeaf{232}allgemein in dem Sinne, dass $M$ die im inneren Gebiet enthaltene Masse bezeichnet. Dabei sind Massen, die etwa auf der Oberfläche selbst stetig verteilt sind, zu den inneren Massen zu zählen oder von ihnen auszuschließen, je nachdem, ob man für $\frac{\mathrm dV}{\mathrm dp}$ den auf die Außenseite oder den auf die Innenseite bezogenen Wert gewählt hat.

Sind demnach im Inneren des Gebiets gar keine Massen enthalten und versteht man unter $\frac{\mathrm dV}{\mathrm dp}$ stets den auf die Innenseite bezogenen Wert, so gilt
\[
\int\frac{\mathrm dV}{\mathrm dp}.\mathrm ds=0
\]
\Article{24}
Wir gehen von denselben Voraussetzungen aus wie am Ende des vorhergehenden Artikels und bezeichnen das betrachtete Gebiet mit $T$. In seinem Element $\mathrm dT$ bezeichnen wir die gesamte Kraft, die von Massen außerhalb des Gebiets oder von stetig auf der Oberfläche verteilten Massen herrührt, mit $q$. Damit erhalten wir den folgenden wichtigen Lehrsatz.

\textsc{Lehrsatz.} Es gilt
\[
\int V\frac{\mathrm dV}{\mathrm dp}.\mathrm ds=-\int qq\,\mathrm dT
\]
wenn das erste Integral über die ganze Fläche und das zweite über das ganze Gebiet $T$ erstreckt wird.

\textit{Beweis.} Wir führen rechtwinklige Koordinaten $x,y,z$ ein und betrachten zunächst eine zur $x$-Achse parallele Gerade, die das Gebiet $T$ schneidet. Dabei haben $y,z$ konstante Werte. Aus der Identität
\[
\frac{\mathrm d}{\mathrm dx}(V\frac{\mathrm dV}{\mathrm dx})=(\frac{\mathrm dV}{\mathrm dx})^2+V\frac{\mathrm d\mathrm dV}{\mathrm dx^2}
\]
folgt, dass das Integral
\[
\int((\frac{\mathrm dV}{\mathrm dx})^2+V\frac{\mathrm d\mathrm dV}{\mathrm dx^2})\,\mathrm dx
\]
über das innerhalb von $T$ liegende Stück der Geraden gleich der Differenz der beiden Werte von $V\frac{\mathrm dV}{\mathrm dx}$ am Anfangs- und Endpunkt ist, sofern die Gerade die Grenzfläche nur zweimal schneidet. Allgemein ist es $=\sum\varepsilon V\frac{\mathrm dV}{\mathrm dx}$, wobei für $V\frac{\mathrm dV}{\mathrm dx}$ die einzelnen Werte an den ver\SourceLeaf{233}schiedenen Schnittpunkten einzusetzen sind und $\varepsilon$ an den ungeraden Schnittpunkten (am ersten, dritten usw.) $=-1$ und an den geraden $=+1$ ist. Betrachten wir weiter längs dieser Geraden das prismatische Gebiet, dessen Querschnitt das Rechteck $\mathrm dy.\mathrm dz$ und dessen Element somit $\mathrm dx.\mathrm dy.\mathrm dz$ ist, so ist das Integral
\[
\int((\frac{\mathrm dV}{\mathrm dx})^2+V\frac{\mathrm d\mathrm dV}{\mathrm dx^2})\,\mathrm dT
\]
über den in diesem Prisma liegenden Teil von $T$ $=\sum\varepsilon V\frac{\mathrm dV}{\mathrm dx}.\mathrm dy.\mathrm dz$. Dieses Prisma schneidet aus der Grenzfläche zwei oder allgemein eine gerade Anzahl von Stücken aus. Bezeichnen wir jedes dieser Stücke mit $\mathrm ds$ und mit $\xi$ den Winkel zwischen der $x$-Achse und der nach innen gerichteten Normalen auf $\mathrm ds$, so gilt $\mathrm dy.\mathrm dz=\pm\cos\xi.\mathrm ds$, wobei das obere Zeichen für die ungeraden und das untere für die geraden Schnittpunkte genommen wird. Das obige Integral wird folglich
\[
=-\sum V\frac{\mathrm dV}{\mathrm dx}.\cos\xi.\mathrm ds
\]
wobei sich die Summe über sämtliche betreffenden Flächenelemente erstreckt. Zerlegt man nun das ganze Gebiet $T$ in solche prismatischen Elemente, so füllen auch alle zugehörigen Stücke der Fläche diese ganz aus. Somit gilt
\[
\int((\frac{\mathrm dV}{\mathrm dx})^2+V\frac{\mathrm d\mathrm dV}{\mathrm dx^2})\,\mathrm dT=-\int V\frac{\mathrm dV}{\mathrm dx}.\cos\xi.\mathrm ds
\]
wobei sich die erste Integration über das ganze Gebiet $T$ und die zweite über die ganze Fläche erstreckt. Offenbar ist nun $\cos\xi$ gleich der partiellen Ableitung $\frac{\mathrm dx}{\mathrm dp}$, denn $p$ hat die in Artikel 23 festgelegte Bedeutung, und $x$ lässt sich als Funktion von $p$ und zwei weiteren veränderlichen Größen auffassen, die die einzelnen Punkte der Fläche voneinander unterscheiden. Folglich gilt
\[
\int((\frac{\mathrm dV}{\mathrm dx})^2+V\frac{\mathrm d\mathrm dV}{\mathrm dx^2}).\mathrm dT=-\int V\frac{\mathrm dV}{\mathrm dx}.\frac{\mathrm dx}{\mathrm dp}.\mathrm ds
\]
Es ist außerdem unmittelbar klar, dass, wenn die Fläche selbst Massen enthält und $\frac{\mathrm dV}{\mathrm dx}$ somit zwei verschiedene Werte hat, hier immer der auf das innere Gebiet bezogene Wert gemeint ist.

Mit ganz ähnlichen Überlegungen findet man
\SourceLeaf{234}\[
\begin{gathered}
\int((\frac{\mathrm dV}{\mathrm dy})^2+V\frac{\mathrm d\mathrm dV}{\mathrm dy^2})\,\mathrm dT=-\int V\frac{\mathrm dV}{\mathrm dy}.\frac{\mathrm dy}{\mathrm dp}.\mathrm ds\\
\int((\frac{\mathrm dV}{\mathrm dz})^2+V\frac{\mathrm d\mathrm dV}{\mathrm dz^2})\,\mathrm dT=-\int V\frac{\mathrm dV}{\mathrm dz}.\frac{\mathrm dz}{\mathrm dp}.\mathrm ds
\end{gathered}
\]

Addiert man nun diese drei Gleichungen und beachtet, dass im Gebiet $T$
\[
\begin{gathered}
\frac{\mathrm d\mathrm dV}{\mathrm dx^2}+\frac{\mathrm d\mathrm dV}{\mathrm dy^2}+\frac{\mathrm d\mathrm dV}{\mathrm dz^2}=0\\
(\frac{\mathrm dV}{\mathrm dx})^2+(\frac{\mathrm dV}{\mathrm dy})^2+(\frac{\mathrm dV}{\mathrm dz})^2=qq
\end{gathered}
\]
und an der Grenzfläche
\[
\frac{\mathrm dV}{\mathrm dx}.\frac{\mathrm dx}{\mathrm dp}+\frac{\mathrm dV}{\mathrm dy}.\frac{\mathrm dy}{\mathrm dp}+\frac{\mathrm dV}{\mathrm dz}.\frac{\mathrm dz}{\mathrm dp}=\frac{\mathrm dV}{\mathrm dp}
\]
gilt, so erhält man $\int qq\,\mathrm dT=-\int V.\frac{\mathrm dV}{\mathrm dp}.\mathrm ds$, also unseren Lehrsatz selbst. Zieht man auch den letzten Satz des vorhergehenden Artikels heran, lässt er sich noch allgemeiner so ausdrücken
\[
\int qq\,\mathrm dT=\int(A-V)\frac{\mathrm dV}{\mathrm dp}.\mathrm ds
\]
wobei $A$ eine beliebige Konstante bezeichnet.

\Article{25}
\textsc{Lehrsatz.} Hat das Potential $V$ unter denselben Voraussetzungen wie im vorhergehenden Artikel an allen Punkten der Grenzfläche des Gebiets $T$ denselben Wert, so gilt dieser Wert auch für sämtliche. Punkte des Gebiets selbst, und die Kräfte heben sich im ganzen Gebiet vollständig auf.

\textit{Beweis.} Setzt man im erweiterten Lehrsatz des vorhergehenden Artikels für $A$ den konstanten Randwert des Potentials ein, so gilt offenbar $\int qq\,\mathrm dT=0$. Daher gilt notwendig $q=0$ an jedem Punkt des Gebiets $T$ und somit auch $\frac{\mathrm dV}{\mathrm dx}=0$, $\frac{\mathrm dV}{\mathrm dy}=0$, $\frac{\mathrm dV}{\mathrm dz}=0$. Folglich ist $V$ im ganzen Gebiet $T$ konstant.

\Article{26}
\textsc{Lehrsatz.} Liegen Massen ausschließlich innerhalb des endlichen Gebiets $T$ oder sind sie auch ganz oder teilweise stetig auf des%
\SourceLeaf{235}%
sen Oberfläche $S$ verteilt und hat ihr Potential in allen Punkten von $S$ einen konstanten Wert $=A$, so ist das Potential in jedem Punkt $O$ des unendlichen äußeren Gebiets $T'$

\noindent erstens, wenn $A=0$ gilt, ebenfalls $=0$,

\noindent zweitens, wenn $A$ nicht $=0$ ist, kleiner als $A$ und hat dasselbe Vorzeichen wie $A$.

\textit{Beweis.} I. Zunächst ist zu zeigen, dass das Potential in $O$ keinen Wert außerhalb der Grenzen $0$ und $A$ annehmen kann. Nehmen wir an, das Potential habe in $O$ einen solchen Wert $B$, und bezeichnen wir mit $C$ eine beliebige Größe, die sowohl zwischen $B$ und $0$ als auch zwischen $B$ und $A$ liegt. Lässt man von $O$ in alle Richtungen gerade Linien ausgehen, so gibt es auf jeder einen Punkt $O'$, in dem das Potential $=C$ wird, wobei die ganze Linie $OO'$ im Gebiet $T'$ liegt. Das folgt unmittelbar daraus, dass sich das Potential stetig ändert. Setzt man die gerade Linie hinreichend weit fort, so geht es entweder von $B$ in $A$ über oder nimmt unendlich ab, je nachdem, ob die Linie die Fläche $S$ trifft oder nicht (siehe die Bemerkung am Ende des 21.\ Artikels). Die Gesamtheit aller Punkte $O'$ bildet dann eine geschlossene Fläche. Da das Potential auf ihr konstant $=C$ ist, muss es nach dem Lehrsatz des vorhergehenden Artikels in allen Punkten des von dieser Fläche eingeschlossenen Gebiets denselben Wert haben. In $O$ hat es aber den von $C$ verschiedenen Wert $B$. Die Annahme führt daher notwendig zu einem Widerspruch.

Für den Fall $A=0$ ist unser Lehrsatz damit vollständig bewiesen. Für den zweiten Fall, in dem $A$ nicht $=0$ ist, folgt daraus, dass das Potential in keinem Punkt von $T'$ größer als $A$ sein oder das entgegengesetzte Vorzeichen haben kann.

II. Um unseren Beweis für den zweiten Fall zu vervollständigen, beschreiben wir um $O$ als Mittelpunkt eine Kugeloberfläche mit Radius $R$, der kleiner ist als der geringste Abstand des Punkts $O$ von $S$. Wir zerlegen die Oberfläche in Elemente $\mathrm ds$ und bezeichnen das Potential in jedem Element mit $V$. Das Potential in $O$ bezeichnen wir wieder mit $B$. Nach dem Lehrsatz des 20.\ Artikels ergibt sich dann für das Integral über die ganze Kugeloberfläche
\[
\int V\mathrm ds=4\pi RRB,\quad\text{und folglich}\quad\int(V-B)\mathrm ds=0
\]
\SourceLeaf{236}Diese Gleichheit kann aber nur gelten, wenn entweder $V$ in allen Punkten der Kugeloberfläche konstant $=B$ ist oder $V$ in verschiedenen Teilen der Kugeloberfläche in entgegengesetztem Sinn von $B$ abweicht. Im ersten Fall müsste das Potential nach Art.\ 25 im ganzen Inneren der Kugel und daher nach Art 21 im gesamten unendlichen Gebiet $T'$ konstant und $=0$ sein. Das widerspricht der Annahme, dass es am Rand dieses Gebiets, auf der Fläche $S$, von $0$ verschieden ist, denn von dort aus kann es sich nicht sprunghaft ändern. Der zweite Fall hingegen würde dem unter I: Bewiesenen widersprechen, wenn $B$ entweder $=0$ oder $=A$ wäre. Daher muss $B$ notwendigerweise \textit{zwischen} $0$ und $A$ liegen.

\Article{27}
\textsc{Lehrsatz.} Im Lehrsatz des vorhergehenden Artikels kann der erste Fall, also der Wert $0$ für das konstante Potential $A$, nur eintreten, wenn die Summe aller Massen selbst $=0$ ist. Der zweite Fall kann nur eintreten, wenn diese Summe nicht $=0$ ist.

\textit{Beweis.} Sei $\mathrm ds$ ein Element der Oberfläche einer beliebigen Kugel, die das Gebiet $T$ einschließt, $R$ ihr Radius, $M$ die Summe aller Massen und $V$ ihr Potential in $\mathrm ds$. Nach dem Lehrsatz des 20.\ Artikels gilt für das Integral $\int V\mathrm ds=4\pi RM$. Nach dem vorhergehenden Lehrsatz ist im ersten Fall, also für $A=0$, das Potential $V$ in allen Punkten der Kugeloberfläche $=0$. Im zweiten Fall ist es hingegen kleiner als $A$ und hat dasselbe Vorzeichen. Somit gilt im ersten Fall $4\pi RM=0$, also $M=0$. Im zweiten Fall hingegen müssen $4\pi RM$ und damit auch $M$ dasselbe Vorzeichen wie $A$ haben. Zugleich ergibt sich, dass in diesem zweiten Fall $4\pi RM$ kleiner ist als $\int A\mathrm ds$ oder $4\pi RRA$, mithin $M$ kleiner als $RA$ oder $A$ größer als $\frac{M}{R}$.

Der zweite Teil dieses Lehrsatzes lässt sich zusammen mit dem Lehrsatz des vorhergehenden Artikels offenbar auch folgendermaßen ausdrücken.

Sind Massen in einem von einer geschlossenen Fläche begrenzten Gebiet enthalten oder auch teilweise stetig auf der Fläche selbst verteilt, ist ihre algebraische Summe $=0$ und hat ihr Potential in allen Punkten der Fläche einen konstanten Wert, so muss dieser Wert selbst $=0$
\SourceLeaf{237}sein und zugleich im gesamten unendlichen äußeren Gebiet gelten. Folglich hebt sich die Wirkung der von diesen Massen ausgehenden Kräfte im ganzen äußeren Gebiet vollständig auf.

\Article{28}
Man sieht leicht, dass sämtliche Schlüsse der beiden vorhergehenden Artikel gültig bleiben, wenn $S$ eine nicht geschlossene Fläche ist und die Massen ausschließlich auf ihr liegen. Hier fällt das Gebiet $T$ ganz weg. Alle Punkte, die nicht auf der Fläche selbst liegen, gehören zum unendlichen äußeren Gebiet. Hat das Potential auf der Fläche überall einen konstanten, von $0$ verschiedenen Wert $A$, so hat es außerhalb der Fläche überall einen kleineren Wert mit demselben Vorzeichen.

Die Aussage zum ersten Fall, $A=0$, gilt auch hier, ist aber inhaltsleer. Denn in diesem Fall ist das Potential $V$ in allen Punkten des Raumes $=0$. Daher gilt auch überall $\frac{\mathrm dV}{\mathrm dt}=0$, wenn $t$ eine beliebige gerade Linie bezeichnet. Daraus folgt nach Art.\ 18 leicht, dass die Dichte auf der Fläche überall $=0$ sein muss, die Fläche also überhaupt keine Masse enthalten kann.

Diese letzte Bemerkung gilt übrigens allgemein, wenn die Massen nur auf der Fläche selbst liegen sollen, auch wenn sie geschlossen ist. Denn nach dem Lehrsatz des 25.\ Artikels ist der Wert des Potentials in diesem Fall offenbar auch im ganzen inneren Gebiet $=0$.

\Article{29}
Bevor wir zu den folgenden Untersuchungen übergehen, in denen stetig auf einer Fläche verteilte Massen eine Hauptrolle spielen, ist ein wesentlicher Unterschied zwischen den Verteilungen hervorzuheben. Entweder werden nur Massen mit einem einzigen Vorzeichen zugelassen (die wir der Kürze wegen immer als positiv betrachten werden), oder auch Massen mit entgegengesetzten Vorzeichen. Ist eine Masse $M$ so auf einer Fläche verteilt, dass auf jedes Flächenelement $\mathrm ds$ die Masse $m\mathrm ds$ entfällt, so heißt $m$ wie bisher die Dichte, und $\int m\mathrm ds$ über die ganze Fläche erstreckt ist $=M$. Wir nennen dies eine \textit{gleichartige} Verteilung, wenn $m$ überall positiv oder wenigstens nirgends negativ ist. Ist dagegen $m$ an einigen Stellen positiv, an anderen negativ, so soll die Vertei%
\SourceLeaf{238}%
lung \textit{ungleichartig} heißen. Dabei ist $M$ nur die algebraische Summe der Massenanteile oder der absolute Unterschied zwischen den positiven und den negativen Massen. Ein ganz besonderer Fall einer ungleichartigen Verteilung liegt vor, wenn $M=0$. Dann mag es freilich befremdlich erscheinen, weiterhin davon zu sprechen, die Masse $0$ sei über die Fläche verteilt.

\Article{30}
Wie auch immer eine Masse $M$ \textit{gleichartig} über eine Fläche verteilt sein mag, so ist doch klar, dass das daraus entstehende, überall positive Potential $V$ in jedem Punkt der Fläche größer als $\frac{M}{r}$ ist, wobei $r$ den größten Abstand zweier Punkte der Fläche bezeichnet. Genau diesen Wert könnte das Potential nur an einem Endpunkt der Strecke $r$ annehmen, wenn die gesamte Masse im anderen Endpunkt konzentriert wäre. Dieser Fall kommt hier gar nicht in Betracht, da nur von stetiger Verteilung die Rede sein soll, bei der jedem Flächenelement $\mathrm ds$ nur eine unendlich kleine Masse $m\mathrm ds$ entspricht. Das über die ganze Fläche erstreckte Integral $\int Vm\mathrm ds$ ist also in jedem Fall größer als $\int\frac{M}{r}m\mathrm ds$ oder $\frac{MM}{r}$, und daher muss es eine gleichartige Verteilung geben, für die dieses Integral einen Minimalwert annimmt. Hier lässt sich vorab eines der Ziele der folgenden Untersuchungen nennen. Es soll bewiesen werden, dass bei einer solchen Verteilung, bei der $\int Vm\mathrm ds$ seinen Minimalwert annimmt, das Potential $V$ in jedem Punkt der Fläche denselben Wert hat, dass dabei kein Teil der Fläche unbelegt bleiben kann und dass es nur eine einzige solche Verteilung gibt. Der Kürze wegen wollen wir die Untersuchung aber von Anfang an in einer allgemeineren Form durchführen.

\Article{31}
Es sei $U$ eine Größe, die in jedem Punkt der Fläche einen bestimmten endlichen Wert hat, der sich stetig ändert. Das Integral
\[
\Omega=\int(V-2U)m\mathrm ds
\]
über die ganze Fläche kann dann je nach der gleichartigen Verteilung der Masse $M$ sehr unterschiedliche Werte haben. Offenbar
\SourceLeaf{239}muss es jedoch eine solche Verteilung geben, für die dieses Integral einen Minimalwert annimmt. Nun soll der folgende

\textsc{Lehrsatz} bewiesen werden. Bei einer solchen Verteilung

1. hat die Differenz $V-U=W$ überall auf der Fläche, wo diese tatsächlich mit Anteilen von $M$ belegt ist, einen konstanten Wert,

2. muss, wenn Teile der Fläche unbelegt bleiben, $W$ dort größer als jener konstante Wert sein oder kann wenigstens nicht kleiner sein.

I. Zunächst ist zu beweisen, dass, wenn man eine Verteilung durch eine andere ersetzt, die sich unendlich wenig von ihr unterscheidet, indem man $m+\mu$ an die Stelle von $m$ setzt, die daraus entstehende Variation von $\Omega$ durch $2\int W\mu\mathrm ds$ ausgedrückt wird.

Bezeichnen wir die Variationen von $\Omega$ und $V$ mit $\delta\Omega$ und $\delta V$, so gilt nämlich
\[
\delta\Omega=\int\delta V.m\mathrm ds+\int(V-2U)\mu\mathrm ds
\]
Zugleich gilt aber $\int\delta V.m\mathrm ds=\int V\mu\mathrm ds$, wie leicht aus dem Lehrsatz des 19.\ Artikels folgt. Denn $\delta V$ ist das Potential derjenigen Massenverteilung, bei der $\mu$ die Dichte in jedem Flächenelement angibt. Was hier $V,m,\delta V,\mu$ ist, lässt sich also dort für $V,K,v,k$ einsetzen, ebenso $\mathrm ds$ sowohl für $\mathrm dS$ als auch für $\mathrm ds$. Folglich gilt
\[
\delta\Omega=\int(2V-2U)\mu\mathrm ds=2\int W\mu\mathrm ds
\]

II. Die Variationen $\mu$ müssen offenbar allgemein die Bedingung $\int\mu\mathrm ds=0$ erfüllen. Für die vorliegende Untersuchung gilt außerdem eine zweite Bedingung. In den unbelegten Teilen der Fläche, falls solche vorhanden sind, darf $\mu$ nicht negativ sein, da die Verteilung sonst aufhören würde, gleichartig zu sein.

III. Nehmen wir nun an, dass bei einer bestimmten Verteilung von $M$ die Größe $W$ in verschiedenen Teilen der Fläche ungleiche Werte hat. Sei $A$ eine Größe zwischen den ungleichen Werten von $W$. Sei $P$ das Stück der Fläche, in dem die Werte von $W$ größer, und $Q$ dasjenige, in dem sie kleiner sind als $A$. Ferner seien $p,q$ gleich große Stücke der Fläche, das erste zu $P$ und das zweite zu $Q$ gehörig. Unter diesen Voraussetzungen geben wir der Variation
\SourceLeaf{240}von $m$ überall in $p$ den konstanten negativen Wert $\mu=-\nu$, überall in $q$ hingegen den positiven Wert $\mu=\nu$ und in allen übrigen Teilen der Fläche den Wert $0$. Damit ist offenbar die erste Bedingung in II erfüllt. Die zweite verlangt außerdem, dass $p$ keine unbelegten Teile enthält. Das lässt sich immer erreichen, sofern nicht das ganze Stück $P$ unbelegt ist.

Das Ergebnis ist jedoch, dass $\delta\Omega$ einen negativen Wert erhält. Dies sieht man leicht, wenn man diese Variation in der Form $2\int(W-A)\mu\mathrm ds$ schreibt.

Daraus folgt, dass bei einer gegebenen Verteilung, wenn entweder im belegten Teil der Fläche ungleiche Werte von $W$ auftreten oder bei gleichen Werten im belegten Teil kleinere Werte im unbelegten Teil vorkommen, eine Änderung der Verteilung eine Verringerung von $\Omega$ ermöglicht. Beim Minimalwert müssen daher die im obigen Lehrsatz ausgesprochenen Bedingungen erfüllt sein.

\Article{32}
Verbinden wir nun in unserem besonderen Fall (Art.\ 30), in dem $U=0$ gilt, also $W$ das bloße Potential der auf die Fläche verteilten Masse und $\Omega$ das Integral $\int Vm\mathrm ds$ bezeichnet, den Lehrsatz des vorhergehenden Artikels mit dem im 28.\ Artikel angeführten, so folgt unmittelbar, dass beim Minimalwert von $\int Vm\mathrm ds$ die Fläche überhaupt keine unbelegten Teile haben kann. Andernfalls wäre, selbst wenn die ganze Fläche geschlossen ist, der belegte Teil als nicht geschlossene Fläche und der unbelegte Teil ihr gegenüber als zum äußeren Gebiet gehörig zu betrachten. Dort müsste also nach Art.\ 28 das Potential einen kleineren Wert als auf der belegten Fläche haben, während der Lehrsatz des vorhergehenden Artikels einen kleineren Wert ausschließt.

Damit ist bewiesen, dass es eine gleichartige Verteilung einer gegebenen Masse über die ganze Fläche gibt, bei der kein Teil leer bleibt und die in allen Punkten der Fläche dasselbe Potential hervorbringt. Zum vollständigen Beweis des im 30.\ Artikel aufgestellten Lehrsatzes fehlt nun noch der Nachweis, dass es nur eine einzige solche Verteilung geben kann. Dieser wird weiter unten als Teil eines allgemeineren Lehrsatzes folgen.

Dass kein Teil der Fläche unbelegt bleiben darf, wenn $\int Vm\mathrm ds$ seinen Minimalwert annehmen soll, lässt sich offenbar auch folgendermaßen ausdrücken.
\SourceLeaf{241}Bei jeder Verteilung, bei der ein endliches Stück der Fläche leer bleibt, übersteigt das Integral $\int Vm\mathrm ds$ seinen Minimalwert um eine endliche Differenz.

\Article{33}
Der eigentliche Lebensnerv der im 31.\ Artikel entwickelten Beweisführung liegt darin, dass die Existenz eines Minimalwerts für $\Omega$ unmittelbar einleuchtet, solange man sich auf die gleichartigen Verteilungen einer gegebenen Masse beschränkt. Wäre sie ohne diese Beschränkung ebenso einleuchtend, so würden die dortigen Schlüsse unmittelbar zu dem Ergebnis führen, \textit{dass es stets, wenn nicht eine gleichartige, so doch eine ungleichartige Verteilung der gegebenen Masse gibt, bei der $W=V-U$ in allen Punkten der Fläche einen konstanten Wert hat}, da dann die zweite Bedingung (Art.\ 31.\ II) entfällt. Sobald wir aber die Beschränkung auf gleichartige Verteilungen aufgeben, geht jene unmittelbare Einsicht verloren. Deshalb müssen wir den strengen Beweis dieses wichtigsten Satzes unserer gesamten Untersuchung auf einem etwas künstlicheren Weg suchen. Der folgende scheint am einfachsten zum Ziel zu führen.

Wir betrachten zunächst drei verschiedene Massenverteilungen, bei denen wir anstelle der allgemeinen Zeichen für die Dichte $m$ und das Potential $V$ die folgenden besonderen Zeichen verwenden.
\[
\begin{array}{rll}
\mathrm{I.}&m=m^0,&V=V^0\\
\mathrm{II.}&m=m',&V=V'\\
\mathrm{III.}&m=\mu,&V=v
\end{array}
\]
Verteilung I ist diejenige gleichartige Verteilung der positiven Masse $M$, für die $\int Vm\mathrm ds$ seinen Minimalwert annimmt.

II. ist die gleichartige Verteilung derselben Masse $M$, für die $\int(V-2\varepsilon U)m\mathrm ds$ seinen Minimalwert annimmt, wobei $\varepsilon$ einen beliebigen konstanten Koeffizienten bezeichnet.

III. hängt so von I und II ab, dass $\mu=\frac{m'-m^0}{\varepsilon}$, und ist alo eine ungleichartige Verteilung, in der die Gesamtmasse $=0$ ist.

Nach dem im 31.\ Artikel Bewiesenen ist nun $V^0$ konstant auf der ganzen
\SourceLeaf{242}Fläche; $V'-\varepsilon U$ ist konstant auf der Fläche, soweit sie bei der zweiten Verteilung belegt ist, und daher auf demselben Stück der Fläche auch $v-U$, denn $v=\frac{V'-V^0}{\varepsilon}$.

Ob bei der zweiten Verteilung die ganze Fläche belegt ist oder ein größeres oder kleineres Stück unbelegt bleibt, hängt vom Koeffizienten $\varepsilon$ ab. Da die zweite Verteilung in die erste übergeht, wenn $\varepsilon=0$ gilt, wird im Allgemeinen das für einen bestimmten Wert von $\varepsilon$ unbelegt gebliebene Stück der Fläche kleiner werden, wenn $\varepsilon$ abnimmt, und sich schon ganz füllen, bevor $\varepsilon$ den Wert $0$ erreicht. In Ausnahmefällen kann es sich jedoch auch so verhalten, dass stets ein Stück unbelegt bleibt, solange $\varepsilon$ von $0$ verschieden ist und nicht das entgegengesetzte Vorzeichen annimmt. Für unseren Zweck genügt es, $\varepsilon$ unendlich klein anzunehmen. Dann lässt sich leicht zeigen, dass jedenfalls kein endliches Flächenstück unbelegt bleiben kann. Andernfalls müsste nach der Schlussbemerkung von Art.\ 32 das Integral $\int V'm'\mathrm ds$ um eine endliche Differenz größer sein als $\int V^0m^0\mathrm ds$. Bezeichnen wir diese Differenz mit $e$, so ist die Differenz der beiden Integrale
\[
\int(V'-2\varepsilon U)m'\mathrm ds-\int(V^0-2\varepsilon U)m^0\mathrm ds=e-2\varepsilon\int U(m'-m^0)\mathrm ds
\]
die für ein unendlich kleines $\varepsilon$ positiv bleibt. Das widerspricht der Voraussetzung, dass $\int(V-2\varepsilon U)m\mathrm ds$ bei der zweiten Verteilung seinen Minimalwert hat.

Daraus folgt, dass, wenn man in der dritten Verteilung für $\mu$ den Grenzwert von $\frac{m'-m^0}{\varepsilon}$ bei unendlicher Abnahme von $\varepsilon$ nimmt, $v-U$ auf der ganzen Fläche einen konstanten Wert hat.

Bilden wir nun eine vierte Verteilung, indem wir $m=m^0+\mu$ setzen, sodass die Gesamtmasse $=M$ bleibt, so ist das daraus entstehende Potential $=V^0+v$. Es übertrifft daher auf der ganzen Fläche die Größe $U$ um die konstante Differenz $V^0+v-U$. Damit ist der oben ausgesprochene Lehrsatz bewiesen.

\Article{34}
Es bleibt noch zu beweisen, dass es nur eine Verteilung einer gegebenen Masse $M$ geben kann, bei der $V-U$ auf der ganzen Fläche konstant ist. Gäbe es zwei verschiedene Verteilungsar%
\SourceLeaf{243}%
ten mit dieser Eigenschaft und bezeichneten wir $m$ und $V$ bei der ersten mit $m'$, $V'$, bei der zweiten mit $m''$, $V''$, so hätte eine dritte Massenverteilung, bei der $m=m'-m''$ gesetzt wird, das Potential $=V'-V''$, das somit konstant wäre, und die Gesamtmasse $=0$. Dieses konstante Potential müsste daher nach Art.\ 27 notwendig $=0$ sein. Folglich müsste nach Art.\ 28 auch $m'-m''=0$ gelten, die beiden Verteilungen wären also identisch.

Schließlich ist noch zu erwähnen, dass es stets eine Massenverteilung gibt, bei der die Differenz $V-U$ einen \textit{gegebenen} konstanten Wert annimmt. Bezeichnet $\alpha$ einen beliebigen konstanten Koeffizienten und behalten wir die Bezeichnungen für die erste und dritte Verteilung des vorhergehenden Artikels bei, so ist das Potential der Verteilung, bei der $m=\alpha m^0+\mu$ gesetzt wird, $=\alpha V^0+v$. Der konstanten Differenz $\alpha V^0+v-U$ lässt sich durch geeignete Wahl des Koeffizienten $\alpha$ jeder beliebige Wert geben. Die Gesamtmasse dieser Verteilung ist dann aber nicht mehr beliebig, sondern $=\alpha M$. Auf dieselbe Weise wie zuvor sieht man, dass auch diese Verteilungsbedingung nur auf eine einzige Art erfüllt werden kann.

\Article{35}
Die tatsächliche Bestimmung der Massenverteilung auf einer gegebenen Fläche für jede vorgeschriebene Form von $U$ übersteigt in den meisten Fällen das, was die Analysis auf ihrem gegenwärtigen Stand leisten kann. Der einfachste Fall, in dem sie uns möglich ist, ist eine ganze Kugeloberfläche. Wir wollen jedoch gleich den allgemeineren Fall behandeln, in dem die Fläche nur sehr wenig von einer Kugeloberfläche abweicht und Größen höherer Ordnung als die Abweichung selbst vernachlässigt werden dürfen.

Es sei $R$ der Radius der Kugel, $r$ der Abstand eines beliebigen Punkts im Raum von ihrem Mittelpunkt, $u$ der Winkel zwischen $r$ und einer festen Geraden und $\lambda$ der Winkel zwischen der durch diese Gerade und $r$ gelegten Ebene und einer festen Ebene. Der Abstand eines beliebigen Punkts auf der gegebenen geschlossenen Fläche vom Mittelpunkt der Kugel sei $=R(1+\gamma z)$, wobei $\gamma$ ein sehr kleiner konstanter Faktor ist, dessen höhere Potenzen vernachlässigt werden, während $z$ ebenso wie $U$ eine Funktion von $u$ und $\lambda$ ist.

\SourceLeaf{244}Das Potential $V$ der auf der Kugeloberfläche verteilten Masse wird in jedem Punkt des äußeren Gebiets durch eine nach Potenzen von $r$ fallende Reihe ausgedrückt, der wir die Form geben
\[
A^0\frac Rr+A'\left(\frac Rr\right)^2+A''\left(\frac Rr\right)^3+\text{u.\,s.\,f.}
\]
in jedem Punkt des inneren Gebiets dagegen durch die steigende Reihe
\[
B^0+B'\frac rR+B''\left(\frac rR\right)^2+B'''\left(\frac rR\right)^3+\text{u.\,s.\,f.}
\]
Die Koeffizienten $A^0,A',A''$ usw. sind Funktionen von $u$ und $\lambda$, die bekannten partiellen Differentialgleichungen genügen, \textit{Allg.\ Th.\ d.\ Erdm.}\ Art.\ 18. Dasselbe gilt für $B^0,B',B''$ usw. Auf der gegebenen Fläche soll nun das Potential einer gegebenen Funktion von $u$ und $\lambda$ gleich sein, nämlich $V=U$, also
\[
\left(\frac rR\right)^{\frac12}V=(1+\gamma z)^{\frac12}U
\]
Nehmen wir daher an, $(1+\gamma z)^{\frac12}U$ sei in eine Reihe
\[
P^0+P'+P''+P'''+\text{u.\,s.\,w.}
\]
entwickelt, deren einzelne Glieder $P^0$, $P'$, $P''$, $P'''$ usw. ebenfalls den genannten Differentialgleichungen genügen. Berücksichtigen wir, dass die beiden obigen Reihen für das Potential bis zur Fläche selbst gültig bleiben müssen, so folgt
\[
\begin{aligned}
&P^0+P'+P''+P'''+\text{u.\,s.\,f.}\\
&\qquad=A^0(1+\gamma z)^{-\frac12}+A'(1+\gamma z)^{-\frac32}+A''(1+\gamma z)^{-\frac52}+\text{u.\,s.\,f.}\\
&\qquad=B^0(1+\gamma z)^{\frac12}+B'(1+\gamma z)^{\frac32}+B''(1+\gamma z)^{\frac52}+\text{u.\,s.\,f.}
\end{aligned}
\]
Daraus folgt bei Vernachlässigung von Größen der Ordnung $\gamma$
\[
P^0+P'+P''+\text{u.\,s.\,f.}=A^0+A'+A''+\text{u.\,s.\,f.}
\]
und damit (da sich eine Funktion von $u$, $\lambda$ nur auf eine Art in eine Reihe entwickeln lässt, deren Glieder den genannten Differentialgleichungen genügen)
\[
P^0=A^0,\qquad P'=A',\qquad P''=A''\quad\text{u.\,s.\,f.}
\]
Ebenso ergibt sich bei Vernachlässigung von Größen der Ordnung $\gamma$
\[
P^0=B^0,\qquad P'=B',\qquad P''=B''\quad\text{u.\,s.\,f.}
\]

\SourceLeaf{245}Setzt man daher (I)
\[
\begin{aligned}
A^0&=P^0+\gamma a^0,& B^0&=P^0-\gamma b^0\\
A'&=P'+\gamma a',& B'&=P'-\gamma b'\\
A''&=P''+\gamma a'',& B''&=P''-\gamma b''\\
A'''&=P'''+\gamma a''',& B'''&=P'''-\gamma b'''\\
&\text{u.\,s.\,f.}
\end{aligned}
\]
wobei offenbar auch $a^0,a',a'',a'''$ usw. sowie $b^0,b',b'',b'''$ usw. den genannten Differentialgleichungen genügen, und setzt diese Werte in die obigen Gleichungen ein, wobei Größen der Ordnung $\gamma\gamma$ vernachlässigt werden, so ergibt sich nach Division durch $\gamma$ bis auf Fehler der Ordnung $\gamma$ genau
\[
\begin{aligned}
a^0+a'+a''+a'''+\text{u.\,s.\,f.}&=\tfrac12 z(P^0+3P'+5P''+7P'''+\text{u.\,s.\,f.})\\
b^0+b'+b''+b'''+\text{u.\,s.\,f.}&=\tfrac12 z(P^0+3P'+5P''+7P'''+\text{u.\,s.\,f.})
\end{aligned}
\]
Somit gilt bis auf Fehler der Ordnung $\gamma$ genau
\[
b^0=a^0,\qquad b'=a',\qquad b''=a''\quad\text{u.\,s.\,f.}
\]
und folglich bis auf Fehler der Ordnung $\gamma\gamma$ genau (II)
\[
B^0=P^0-\gamma a^0,\qquad B'=P'-\gamma a',\qquad B''=P''-\gamma a''\quad\text{u.\,s.\,f.}
\]

Der Differentialquotient $\frac{\mathrm dV}{\mathrm dr}$ hat auf der Fläche selbst zwei verschiedene Werte. Der zu einem negativen $\mathrm dr$ oder zur inneren Seite gehörige Wert übertrifft den auf der äußeren Seite geltenden um $4\pi m\cos\theta$, wenn $m$ die Dichte am Schnittpunkt und $\theta$ den Winkel zwischen $r$ und der Normalen bezeichnet (Art.\ 13, wo $t,A,k^0$ dasselbe bedeuten wie hier $r,\theta,m$). Man findet diese beiden Werte, indem man die im inneren und äußeren Gebiet geltenden Ausdrücke für $V$ nach $r$ differenziert und dann $r=R(1+\gamma z)$ setzt. Der erste ist also
\[
=\frac1R\{B'+2B''(1+\gamma z)+3B'''(1+\gamma z)^2+\text{u.\,s.\,f.}\}
\]
und der zweite
\[
-\frac1R\{A^0(1+\gamma z)^{-2}+2A'(1+\gamma z)^{-3}+3A''(1+\gamma z)^{-4}+\text{u.\,s.\,f.}\}
\]

\SourceLeaf{246}Multiplizieren wir die Differenz mit $R(1+\gamma z)^{\frac32}$, so erhalten wir
\[
\begin{aligned}
&4\pi mR\cos\theta.(1+\gamma z)^{\frac32}\\
&\quad=\phantom{{}+{}}A^0(1+\gamma z)^{-\frac12}+2A'(1+\gamma z)^{-\frac32}+3A''(1+\gamma z)^{-\frac52}+\text{u.\,s.\,f.}\\
&\qquad+B'(1+\gamma z)^{\frac32}+2B''(1+\gamma z)^{\frac52}+3B'''(1+\gamma z)^{\frac72}+\text{u.\,s.\,f.}
\end{aligned}
\]
Setzen wir hier für $A^0,A'$ usw. die Werte aus I und für $B^0,B'$ u.\,s.\,w.\ die Werte aus II ein und lassen Größen der Ordnung $\gamma\gamma$ weg, so ergibt sich
\[
\begin{aligned}
4\pi mR\cos\theta.(1+\gamma z)^{\frac32}&=P^0+3P'+5P''+7P'''+\text{u.\,s.\,f.}\\
&\quad+\gamma(a^0+a'+a''+a'''+\text{u.\,s.\,f.})\\
&\quad-\tfrac12\gamma z(P^0+3P'+5P''+\text{u.\,s.\,f.})
\end{aligned}
\]
und folglich, da die beiden letzten Reihen einander bis auf Größen der Ordnung $\gamma\gamma$ aufheben,
\[
m=\frac{(1+\gamma z)^{-\frac32}}{4\pi R\cos\theta}.(P^0+3P'+5P''+7P'''+\text{u.\,s.\,f.})
\]
Damit ist die Aufgabe gelöst. Statt $(1+\gamma z)^{-\frac32}$ kann man auch $1-\frac32\gamma z$ schreiben und den Divisor $\cos\theta$ weglassen, insofern, wenigstens im Allgemeinen, $\theta$ von der Ordnung $\gamma$ ist und daher $\cos\theta$ von $1$ nur um eine Größe der Ordnung $\gamma\gamma$ abweicht.

Für eine Kugel, bei der $\gamma=0$ gilt, erhält man exakt
\[
m=\tfrac1{4\pi R}(P^0+3P'+5P''+7P'''+\text{u.\,s.\,f.})
\]
wobei $P^0+P'+P''+P'''+\text{u.\,s.\,f.}$ die Entwicklung von $U$ selbst darstellt.

\Article{36}
Die Größe $U$ wurde in den bisherigen Untersuchungen unbestimmt gelassen. Die Anwendung dieser Untersuchungen auf den Fall, in dem für $U$ das Potential eines gegebenen Massensystems genommen wird, führt uns nun zu dem folgenden wichtigen

\textsc{Lehrsatz.} Jede gegebene Massenverteilung $D$, die entweder ausschließlich auf das von einer geschlossenen Fläche $S$ begrenzte innere Gebiet oder ausschließlich auf das äußere Gebiet beschränkt ist, lässt sich durch eine Massenverteilung $E$ allein auf der Fläche selbst ersetzen, sodass die
\SourceLeaf{247}Wirkung von $E$ der Wirkung von $D$ gleich ist, im ersten Fall in allen Punkten des äußeren Gebiets und im zweiten in allen Punkten des inneren Gebiets.

Dazu ist nur erforderlich, dass, wenn man das Potential von $D$ in jedem Punkt von $S$ mit $U$ und das Potential von $E$ mit $V$ bezeichnet, auf der ganzen Fläche im ersten Fall $V-U=0$ gilt, im zweiten dagegen die Differenz nur konstant wird. Denn $-U$ ist das Potential einer Verteilung $D'$, die $D$ entgegengesetzt ist (jeder Massenanteil wird durch einen entgegengesetzten ersetzt). Somit ist $V-U$ das Potential der gleichzeitig bestehenden Verteilungen $D'$ und $E$. Deren Wirkungen heben sich daher im ersten Fall im gesamten äußeren Gebiet, im zweiten im gesamten inneren Gebiet auf (Artt.\ 27 und 25). Die Wirkungen von $D$ und $E$ sind also in den jeweiligen Gebieten gleich. Außerdem ist im ersten Fall die Gesamtmasse in $E$ gleich der Masse in $D$, während sie im zweiten beliebig bleibt.

Der in der \textit{Intensitas vis magneticae}, Art.\ 2, angekündigte und auch in der \textit{Allgemeinen Theorie des Erdmagnetismus} an verschiedenen Stellen angeführte Lehrsatz erscheint nun als besonderer Fall des hier Bewiesenen.

\Article{37}
Obwohl, wie schon im 35.\ Artikel bemerkt, die tatsächliche vollständige Bestimmung der Verteilung $E$ in den meisten Fällen unüberwindliche Schwierigkeiten bietet, gibt es doch einen Fall, in dem sie sich sehr leicht durchführen lässt und der hier besondere Erwähnung verdient. Das ist der Fall, in dem $U$ konstant, also $S$ eine Gleichgewichtsfläche für das gegebene Massensystem $D$ ist. Man sieht leicht, dass hier nur der Fall zu betrachten ist, in dem $D$ im inneren Gebiet liegt und die Gesamtmasse nicht $=0$ ist. Andernfalls gäbe es überhaupt keine Wirkung, die durch eine Massenverteilung auf $S$ ersetzt werden müsste.

Es sei $O$ ein Punkt der Fläche $S$ und $r$ eine Gerade, die dort die Fläche rechtwinklig schneidet und in der Richtung von innen nach außen als wachsend betrachtet wird. Ferner sei $-C$ der Wert des Differentialquotienten $\frac{\mathrm dU}{\mathrm dr}$ in $O$ und $m$ die Dichte, die bei der Massenverteilung $E$ in $O$ vorliegt. Der Differentialquotient $\frac{\mathrm dV}{\mathrm dr}$ hat in $O$ zwei verschiedene Werte. Der auf die äußere Seite bezogene Wert ist,
\SourceLeaf{248}weil auf der Fläche und im gesamten äußeren Gebiet $V=U$ gilt, gleich dem Differentialquotienten $\frac{\mathrm dU}{\mathrm dr}$, also $=-C$. Der auf die innere Seite bezogene Wert ist dagegen $=0$, weil $V$ auf der Fläche und im gesamten inneren Gebiet konstant ist. Da nun der zweite Wert um $4\pi m$ größer als der erste ist, gilt $4\pi m=C$ oder $m=\frac C{4\pi}$. Offenbar ist $C$ die aus der Massenverteilung $D$ entstehende Kraft und hat dasselbe Vorzeichen wie die Gesamtmasse.

